\quellseite{496}

a Dieser Arauna wird der Jebusiter König gewesen sein zu Jerusalem, und hernach zu Gott bekehret, fromm und selig worden, sich des königreichs verziehen um Gottes willen.

\section{Summarium}

Hie ist eine Historia, wie David, ob er gleich seine Sünde erkennet, und ließ es ihm leid sein, dennoch eine zeitliche Strafe darum tragen muß. Und daß Gott die Sünde pflegt mit Hunger, Krieg und Pestilenz zu strafen. Darum soll man in Gottes Furcht wandeln, und sich für Sünden bewahren, auch Gott ernstlich bitten, daß er für Sünden und Bösheit wolle, so wir anders dergleichen Strafe entfliehen wollen. Unter welchem doch die allerleichteste ist, wenn Gott mit Pestilenz strafet. Denn der Hunger ist zu jämmerlich, so ist der Krieg zu greulich, als der alles mit sich reißet und verwüstet.

\section{Votum}

Der treue Gott erhalte uns in seiner Furcht, daß wir nicht in Sicherheit und andere Sünden gerathen, und strafe ja nicht unsere bisher begangene Sünde nach unserm Verdienst, und behüte uns sonderlich für Krieg und Blutvergießen, aus Barmherzigkeit, Amen.

\part{Das erste Buch von den Königen}

\chapter{Das 1. Kapitel}

\section{Vorrede}

Es ist gewiß und wahr, daß Gamaliel, der fromme Pharisäer sagt in den Geschichten der Apostel im 5. Kapitel. V. 39.: Was aus Gott ist, das kann niemand dämpfen. Denn sind die Menschen klug und stark; Gott ist weit verständiger und mächtiger, und was er will, das hat Fortgang wider aller Welt, ja auch der Höllen Pforten, Hinderniß. Derwegen sehe ein jeder in seinem Beruf auf Gott, und demselben befehle er desselben Anfang und Ausgang, und lasse sich anderer Leute Widerwärtigkeit und des Teufels Hinderniß nicht irren, lustig machen und aufhalten. Gott ist groß von Rath und mächtig von That. Jer. 32, 19. Und wem Er wohl will, dem muß alles wohl gelingen, in seinen guten Sachen, rechtmäßigem Handel und christlichem Wandel. Davon höret auch an das 1. Kapitel des ersten Buchs der Könige.

\section{Salomo zum Könige gesalbet und angeweyhet}
\quellseite{497}

\section{Summarium}

Es ist eine Historia, wie Adonia, Davids Sohn, etliche fürnehmste Amtleute an sich hänget, und ohne Davids Wissen und Willen das Königreich begehrt an sich zu bringen, V. 5. Aber es gehet, wie der Engel sagt Dan. 4, 14.: Daß der Höchste Gewalt hat über der Menschen Königreiche, und gibt sie, wem er will. Adonia wollte König sein, und ward es nicht; Salomon wollte es nicht sein und ward es. Also muß Gottes Rath für sich gehen; aber Menschen Anschläge müssen zurück gehen, sonderlich in solchen Händeln, die Land und Leute betreffen, und die Leute sich bedünken lassen, sie wollen mit eigener Klugheit viel ausrichten.
\quellseite{498}

\section{Summarium}

Sie ist eine Historia, wie Salomon ordentlicher Weise, nicht aus eigenem Fürnehmen, von David, seinem Vater, zum Könige erwählet und von dem Propheten Nathan und dem Priester Zadok gesalbet wird. Und ist hie sonderlich zu merken, wie das Volk für den neuen König betet, daß Gott bei ihm und seinem Regiment sey, V. 35. 40. 47. Denn wo weltliche Obrigkeit ihren Leuten zum Besten fürgehen soll, da ist sonderlich vonnöthen, daß Gott mit
\quellseite{499}
Segen dabei sei. Es will aber Gott um seinen Segen gebeten sein.

\section{Votum}

Wir danken GOtt für alle unsere liebe Obrigkeit hohen und niedern Standes; der wolle die ferner erhalten und durch den heiligen Geist regieren, daß alle ihre Rathschläge gelangen zu seinen Ehren und den Unterthanen zu Heil und Wohlfahrt, Amen.

\chapter{Das 2. Kapitel}

\section{Vorrede}

An wohlgezogenen Kindern erleben die Eltern Ehre und auch Freude. Und kann den Eltern keine größere Freude widerfahren, denn da sie sehen, daß ihre Kinder sind lebendige Werkzeuge Gottes, Gefäße der Ehren, der Gnaden und Barmherzigkeit Gottes, die GOtt recht erkennen, ihm zu Ehren dienen, ihn anrufen, rühmen und preisen. Darnach, aus wohlgezogenen Kindern werden auch Leute, die ihrem Nächsten wohl und fruchtbarlich dienen, oder derer Dienst dem gemeinen bürgerlichen Leben heilsam und nützlich ist. Derwegen die Eltern auf nichts mehr Mühe und Fleiß wenden sollen, denn daß sie ihre Kinder also unterweisen und ziehen, auf daß sie in diesem Leben an ihren Kindern Freude und Ehre haben, und daß nach ihrem Tode dieselben auch Gott in Heiligkeit und Gerechtigkeit dienen und andern Leuten nütze und dienstlich sein mögen. Dies lehret auch das 2. Kapitel.

\section[Davids letzter Wille und Abschied. Salomos Folge und]{Davids letzter Wille und Abschied. Salomos Folge und ernstliche Regierung}
\quellseite{500}

\section{Summarium}

Hie ist ein sehr schön Exempel, wie David erstlich seinen Sohn Salomo vor seinem Absterben mit Ernst vermahnet, daß er in der Furcht Gottes wandeln und sich in allem nach Gottes Wort halten wolle, auf daß er und seine Nachkommen Glück und Heil mögen haben, V. 3. Derhalben soll man solche Historia und Wort in aller Könige und Fürsten, ja in eines jeden Haus-Vaters Kammern malen und schreiben, daß sie mit solchem Unterricht und treulichen Vermahnen ihren Kindern fleißig vorgehen.

Zum andern ermahnet er ihn auch seines Amts, daß, weil ihm Gott das Schwert in die Hände hatte gegeben, er die Untugend wolle strafen und die bösen Buben bei dem Halse nehmen; V. 5. Denn wo weltliche Obrigkeit die Sünde nicht mit Ernst, wie sichs gebühret, strafen will, da muß Gott strafen und gehet die Strafe gemeiniglich darnach über den ganzen Haufen, das ist, über die Obrigkeit und Unterthanen zugleich.
\quellseite{501}

\section{Summarium}

Hie sind schreckliche Exempel, in welchen wir sehen, wie es mit Adonia, Joab und Simei rc. habe ein Ende genommen: nämlich, daß sie ihrem Verdienste nach endlich sind getödtet worden. Auf daß nicht allein jedermann sich für Aufruhr und Blutvergießen hüte, sondern auch die Obrigkeit lerne, wo sie Glück und Heil will haben, daß sie solche Sünde nicht ungestraft lasse, insonderheit, wo wider Gottes Ordnung Blut ist vergossen, wie Joab gethan hatte.

\section{Votum}

O Du, welcher der rechte Vater ist über alles, was da Kinder heißt im Himmel und auf Erden, helfe gnädiglich, daß auch heute die Eltern ihre Kinder fleißig unterweisen und ziehen; und regiere alle Obrigkeit, daß sie alle Laster strafen, und daß Kinder und Unterthanen in seiner heiligen Furcht und christlichem Gehorsam leben, Amen.

\chapter{Das 3. Kapitel}

\section{Vorrede}

Gleich wie die Obrigkeit selbst Gott anrufen soll, daß er zu ihrer Regierung Glück und Segen geben wolle: also sollen die Unterthanen für ihre Obrigkeit auch fleißig beten. Das lehret der heilige Prophet Jeremias die gefangenen Juden, und schreibt ihnen einen Brief ins babylonische Gefängniß, darinnen er sie, die doch von fremder heidnischer Obrigkeit gefangen gehalten waren, also anredet: Jer. 29 V. 7. So spricht der HErr Zebaoth: Suchet der Stadt Bestes, dahin ich euch habe lassen wegführen, und betet für sie zum HErrn. Denn wenn es ihr wohlgehet, so gehet es euch auch wohl. Und wer das thut, der dienet erstlich damit seiner ordentlichen Obrigkeit besser, denn mit allen Zinsen, Geschoß und Tributen; darnach bessert er das Regiment. Denn solche Gebete sind die allerbesten Anlagen und Landsteurer, welche die Unterthanen zu erlegen schuldig und die dem ganzen Regiment besserlich sind. Darauf höret ferner an das 3. Kapitel.

\section{Salomos Ehe, Gebet und Urtheil}
% Scan 502 omitted: Bible narrative only; no non-biblical apparatus or original structural heading on this scan.
\quellseite{503}

\section{Summarium}

Sie ist eine schöne Historia, welche sonderlich weltliche Obrigkeit wohl merken und bedenken soll, wie Salomo nichts anderes von Gott bittet, denn Weisheit und Verstand, auf daß er dem Regiment wohl fürstehen und jedermann recht thun möge; solches Gebet gefällt Gott wohl, V. 9. Darum wo ein Regent solches mit Ernst bittet, will es Gott reichlich geben. Wie denn Salomo geschah, daß seines gleichen für ihn nicht gewesen, noch nach ihm aufkommen ist.

Zum andern ist hie zu merken, daß Gott Salomo verheißen, wo er sich seinem Vater David nachhalte, wolle er ihm langes Leben geben, V. 14. Daraus folget, daß gottlose Regenten um ihrer Sünden willen oft vor der Zeit sterben und verderben müssen. Darum schaue jedermann, daß er gottesfürchtig sei und recht handele.

\section{Votum}

Der allmächtige Gott, in welches Hand stehen alle Regimente auf Erden, wolle unserer lieben Obrigkeit langes Leben geben und erhalten, und sie regieren, auf daß Gottes Wort ausgebreitet, die Gerechtigkeit befördert und die Bosheit verhindert und gestraft werde, Amen.

\chapter{Das 4. Kapitel}

\section{Vorrede}

Was Gott zu Josua sagt, Jos. 1, 8., das wird allen Regenten gesagt: also aber spricht Gott der Herr: Laß das Buch des Gesetzes nicht von deinem Munde kommen, sondern betrachte es Tag und Nacht, auf daß du haltest und thust allerdinge, nach dem, das darinnen geschrieben steht. Alldieweil wird dir gelingen in allem, das du thust und wirst weislich handeln können. O eine gute nöthige Lehre und tröstliche heilsame Verheißung. So helfe der allmächtige Gott, daß alle Regenten in ihrem Amt nicht ihrem eigenen verderbten Willen, auch nicht bösen Rathgebern, sondern Gottes Wort und Willen folgen, und beende mit reicher Gnade ihr Regiment, daß darinnen ihm eine ewige Kirche gesammelt und gemeiner Friede, mit Zucht und Gerechtigkeit, erhalten werde, zu seinem Lob und zu Nutz den Unterthanen, Amen. Darauf höret auch an das 4. Kapitel.

\section{Salomos Amtleute, Herrlichkeit und Weisheit}
\quellseite{504}

\section{Summarium}

Ist eine Historia, wie Salomo seinen Hof bestellet, und Gott ihm großen Reichthum gegeben, V. 7., und einen herrlichen Namen bei allen Heiden gemacht habe, V. 34. Solches alles dienet dazu, daß man lerne, daß Gott Segen, Glück und Heil den Regenten geben und widerfahren lassen wolle, die sich seines Worts halten und sich befleißigen, daß sie ihrem Amt weißlich und getreulich fürstehen.

\section{Votum}

GOtt sei Dank für alle christliche Obrigkeit und für alle Wohlthaten, so er den Unterthanen durch sie erzeiget hat; der verleihe, daß unsere Regenten alle ferner in ihrem Amt fleißig sein nach GOttes Wort, und glücklich nach GOttes Verheißungen, Amen.

\chapter{Das 5. Kapitel}

\section{Vorrede}

Den gemeinen Land- und Stadtfrieden lobet und liebet jedermann; aber gar wenig sind ihr, die ihn recht und christlich gebrauchen. Derwegen laßt uns Gott herzlich danken, daß er uns so lange Zeit Friede verleihet gnädiglich, und schaue ein jeder ja zu, daß er den Frieden in seinem Beruf gebrauche, Gott zu Ehren, seinem Nächsten zu Nutz, und ihm selber zu gutem Gewissen, auf daß man mit Mißbrauch des Friedens zu Krieg und andern Strafen nicht Ursache gebe. Ach, das merket wohl. Denn Krieg ist der größten Strafen eine, als die da die Religion, weltlich und häuslich Regiment, zerstöret und hinweg nimmt, daß Pestilenz Kinderspiel gegen Krieg zu rechnen ist. Darum auch David lieber Pestilenz wählete, denn Krieg, da er unter dreien Strafen eine wählen sollte; wollte also lieber in Gottes, denn Menschen Hände fallen, 2 Sam. 24, 14. Darauf höret ferner an das 5. Capitel.

\section{Salomos Bund mit Hiram, bei Vorbereitung zum Tempelbau}
\quellseite{505}

\section{Summarium}

Hie ist ein fein Exempel, welches sonderlich die Obrigkeit merken und ihm nachfolgen soll: daß Salomo den Frieden, so Gott in seinem Regiment ihm widerfahren läßt, dahin brauchet und wendet, daß er den rechten Gottesdienst ausrichten und fördern möge, V. 3. Denn dahin soll man nicht allein Frieden, sondern auch allen andern Segen, wie ihn Gott gibt, brauchen, daß Gottes Ehre und Wort gefördert, und die Leute dadurch gebessert werden. Da sonst die Welt den Frieden mißbrauchet, entweder zu eignem Wollust und Sicherheit, oder zu Muthwillen und unnöthigen Rumoren.

\section{Votum}

Gott, unserm HErrn, sei Lob und Dank, daß unsere liebe Obrigkeit den Frieden auch erhält; der verleihe ferner, daß wir unter ihr sein Wort und Frieden haben, und allzeit recht' und christlich gebrauchen, auf daß wir in Friede leben und sterben, Amen.

\chapter{Das 6. Kapitel}

\section{Der Tempel Salomos wird in sieben Jahren aufgebaut}

a Dreißig Ellen hoch) im andern Buch der Chronica Kap. 3, 4. spricht der Text: das Haus sei hundert und zwanzig Ellen hoch gewesen, welches ist von des ganzen Hauses Höhe geredet. Sie aber redet er vom untersten Gemach alleine, welches dreißig Ellen hoch war.

b Die Fenster, so inwendig weit und auswendig enge sind, da kann man nicht wohl noch viel hinein sehen; aber sehr hoch
% Scan 506 omitted: Bible narrative only; no non-biblical apparatus or original structural heading on this scan.
\quellseite{507}

\chapter{Das 7. Kapitel}

\section[Verfertigung der königlichen Häuser und der Gefäße des]{Verfertigung der königlichen Häuser und der Gefäße des Tempels}
\quellseite{508}

\chapter{Das 8. Kapitel}

\section{Vorrede}

Zu einem recht christlichen Leben gehören auch diese Stücke: erstlich, daß ihm niemand aus eigener Andacht Gottesdienst soll erdencken, sondern allein thun, was Gott in seinem Worte befohlen hat. Und das lehret Gott selber Ezech. 20, 18., da er spricht: Ich bin der HErr, euer Gott, nach meinen Geboten sollt ihr leben, und meine Rechte sollt ihr halten. Zum andern soll
\quellseite{509}
ein jeder Gott danken für empfangene Wohlthaten. Das lehret auch Gott, der HErr, und spricht Pf. 50, 15.: Rufe mich an in der Noth, so will ich dich erretten, so sollt du mich preisen. Und zum dritten soll man beten für seine eigene und für der Christenheit gemeine Noth. Das lehret St. Paulus 1 Tim. 2, 1. und sagt: Ich ermahne euch, daß ihr für allen Dingen thut Bitte, Gebet, Fürbitte und Danksagung für alle Menschen. Dergleichen, welcher Mensch Buße thut, der diene also Gott und seinem Nächsten; so wird’s eine rechte, Gott, dem HErrn, eine wohlgefällige, und ihm und seinem Nächsten eine nützliche Buße sein. Davon höret auch an dem ersten Theil des 8. Kapitels.

\section{Der Tempel Salomos wird eingeweihet}



\section{Summarium}

In eine Historia, wie man die Lade des HErrn sammt dem andern Kirchengeräthe in den Tempel gebracht, und da Gott geopfert und gedanket hat. Sie muß man erstlich merken, daß Salomo
\quellseite{510}
solches nicht aus eigener Andacht, sondern aus Befehl Gottes gethan habe. Sonst wäre es Gott kein Dienst noch Wohlgefallen gewesen. Zum andern muß man auch darauf sehen, was es für eine Ursache sei, daß Salomo und alles Volk so froh ist, und Gott so loben und danken; nämlich darum, daß der Gottesdienst an einem gewissen Ort verordnet war, da er allweg sollte bleiben, und die Leute gewiß wüßten, wenn sie an dem Ort beteten, opferten und dergleichen thäten, daß es Gott dienet und seinem Befehl gefolget hieß, und da Gott sich finden lassen und sie erhören wollte.

a Brand) wenn das Getreide verschienen, oder von der Sonne verbrannt ist.
\quellseite{511}

b Fremder) dies Haus soll ein Bethaus sein allen Völkern.

\section{Summarium}

Hie ist ein Gebet, daß Gott ehrlich seine Verheißung, so er David gethan, halten und vollstrecken wollte. Auf daß sein Reich so lange bestehen möge, bis daß Christus käme und das ewige Reich durch sein Evangelium anrichtete, V. 29. u. f. Zum andern, daß Gott sich in diesem Hause gnädig wolle lassen finden, und seines Volks, auch der Heiden Gebet in allen Nöthen und Anliegen erhören, V. 25. Und hie ist sonderlich zu merken, daß Salomo allweg hinzusetzt, daß sie Gott erhören wolle, wenn sie ihre Sünde erkennen und den Namen des HErrn bekennen etc. Denn wo Sünde und Unglück uns nicht drücken, daß wirs fühlen, schreien wir zu Gott um Gnade und Hülfe nicht.

\section{Votum}

Der barmherzige GOtt führe, leite und regiere uns, daß wir in seinen Geboten wandeln; und wie er uns bisher aus Gnaden viel Wohlthaten erzeiget hat, also wolle er uns ferner segnen, und mit uns alle gottselige Christen an Leib und Seele, zeitlich und ewig, Amen.

\section{Der andere Theil des 8. Capitels}

\section{Vorrede}

Gleich wie Salomo für alle seine Unterthanen bittet; also soll täglich ein jeder Christ für den andern bitten. Denn also lehret St. Jacobus Kap. 5, 16. und sagt: Betet für einander. Und derwegen leben und wohnen auch die Menschen beisammen, in Städten und Dörfern, nicht daß sie mit einander fressen und saufen, und andere Wollust treiben sollen; sondern daß einer des andern Noth sehen, ihm rathen und helfen und für ihn bitten soll. Und soll solches thun, nicht allein angelangt, sondern auch ungebeten, der anders ein Christ ist; und soll also sein als der erste Psalm sagt V. 2.: Wie ein Baum, der an die Wasserbäche gepflanzet ist, und seine Frucht bringet zu seiner Zeit. Ein guter Baum aber gibt keine Früchte, wenn sie reif sein, auch ungeschüttelt und ungerüttelt; also thut auch ein rechter Christ Gutes, auch ungebeten, zu seiner Zeit, wenn er siehet, daß sein Nächster seiner Hülfe bedarf, und bittet für ihn. Darauf höret nun an den andern Theil des 8. Capitels.
\quellseite{512}

\section{Summarium}

Ist ein feiner Beschluß des Gebets, welches Salomo für sein Reich und Volk gethan hat, da er begehret, Gott wolle solch Gebet erhören, bei ihnen sein, wie er bei ihren Vätern gewesen, sie nicht verlassen, noch die Hand abziehen, sondern ihre Herzen dahin neigen, daß sie in seinen Wegen wandeln und sich nach seinem Worte halten mögen. Auf die Weise sollte eine christliche Obrigkeit alle Tage für sich und ihre Unterthanen beten; so würde Glück und Heil beim Regiment sein.

Das aber hie stehet, daß Salomo habe das Haus also geweihet, V. 63., ist viel ein ander Ding, denn wenn unsere Papisten Kirchen weihen, mit Kreuz machen, segnen und anderm Affenspiel, welches sie ohne Gottes Befehl selbst erdichtet haben. Denn weihen heißt hie nichts anders, denn daß man da habe angefangen zu opfern, und den rechten Gottesdienst zu üben, als, Gottes Gnade und Wohlthat predigen, ihn daselbst loben, preisen, in aller Noth anrufen etc. Von solchem Weihen wissen die Papisten nichts, denn sie üben es ja nicht, und gehen derweil mit ihrem unnützen, stummen und ohne Gottes Befehl erdichtetem Gaukelwerk um.

\section{Votum}

GOtt, der heilige Geist, der ein Geist der Gnaden und Gebets ist, helfe uns allen, daß wir in diesen betrübten und gefährlichen Läuften fleißig für einander beten, und in allem Gehorsam der heiligen zehen Gebote GOttes leben, und behüte uns für aller Abgötterei und falschen Gottesdiensten, Amen.

\chapter{Das 9. Kapitel}

\section{Vorrede}

Wenn wir zu Gott, dem HErrn, unser Gebet thun, so sollen wir mit allem Fleiß einen Unterschied halten zwischen den Gütern und Gaben, die wir bitten. Denn, wenn wir Gott bitten um geistliche ewige Gaben, die zu unserm Heil und Seligkeit gehören, da stehet die Verheißung gewiß und fest, ohne alle Exception und Bedingung, daß uns solche Gott wahrhaftig will gewähren, so gewiß als er lebet und unser Vater ist. Wenn wir aber bitten um zeitliche, leibliche Gaben, da stehet zwar die Verheißung auch gewiß und fest; aber mit der Bedingung, wo dasjenige, so wir bitten, Gott zu Ehren gelanget und uns gut und selig ist. Derwegen auch die lieben Heiligen in der Schrift, wenn sie um vergängliche, leibliche Gaben bitten, allzeit den Willen Gottes zuvor behalten; wie davon lehret
\quellseite{513}
der heilige Apostel Johannes 1. Epistel 5, V. 14.: Das ist die Freudigkeit, die wir haben zu Gott, daß so wir etwas bitten nach seinem Willen, so höret er uns. Und so wir wissen, daß er uns höret, was wir bitten, so wissen wir, daß wir die Bitte haben, die wir von ihm gebeten haben. Dieß soll auch erinnern das 9. Kapitel.

\section[Der HErr erscheint Salomo zum andern mal. Feste]{Der HErr erscheint Salomo zum andern mal. Feste Städte. Schiffahrt}
\quellseite{514}

\section{Summarium}

Ist eine Historia, wie der HErr dem Salomo erschienen, V. 2., und ihm seine Bitte gewähret habe, V. 3. Daher stimmt es, daß hernach Daniel und andere Heiligen, wenn sie im Lande Babel beteten, ihr Angesicht gegen Jerusalem zum Tempel kehreten. Weil Salomo darum gebeten und Gott hernach solches zugesagt hatte, er wollte alsdann hören und helfen, wo sie an diesem Ort, oder so sie nicht im Lande wären, doch gegen diesen Ort beteten, und ihr Angesicht gegen Jerusalem wendeten. So viel aber die Verheißung das leibliche Königreich betrifft und Salomons Nachkommen, siehet man, daß es mit dem Gedinge ist zugesagt: Gott wollte sie dabei erhalten, so fern sie fromm sind und Gottes Wort halten, V. 4. Darum, da sie böse und ungehorsam worden, sind sie, wie der HErr hie dräuet, gestraft und von solcher Herrlichkeit gestoßen worden.

\section{Votum}

Der treue wahrhaftige GOtt erhöre auch unser gläubiges Gebet, daß wir nach seinem Willen thun, und segne und schütze uns und die Unsern in diesen gefährlichen und kümmerlichen Zeiten, und thue solches um JEsu Christi, unsers HErrn und Heilandes willen, Amen.

\chapter{Das 10. Kapitel}

\section{Vorrede}

St. Paulus in der Epistel an die Römer im 15. Kapitel V. 9. sagt, daß die Heiden Gott loben um der Barmherzigkeit willen. Und David im 117. Psalm V. 1. vermahnet: Lobet den HErrn alle Heiden, und preiset ihn alle Völker. Derwegen, weil wir auch geborne Heiden sind, so danken wir heute abermals Gott herzlich, daß er uns zu seinem Volke angenommen und zum Reich Christi hat kommen, und daß er uns seiner Wohlthaten läßt theilhaftig werden. Der helfe, daß wir als sein Volk alhie auf Erden in wahrem Glauben und Gehorsam, und droben im Himmel, in ewiger Seligkeit und Freude, alle leben und ihn, mit allen Auserwählten, für diese und alle andere seine Wohlthaten, loben und preisen. Davon erinnert auch der erste Theil des 10. Kapitels.

\section[Die Königin aus Reicharabien verwundert sich über Salomos]{Die Königin aus Reicharabien verwundert sich über Salomos Weisheit, Reichthum und Herrlichkeit}
\quellseite{515}

a Ebenholz) ist ein Baum in Indienland.

\section{Summarium}

Ist eine Historia, daß Salomo seiner Weisheit halben allenthalben bekannt und hoch berühmt ist gewesen, V. 1. Solches aber hat dazu gedienet, daß durch solch Gerücht etliche auch aus den Heiden gen Jerusalem kamen, und durch solche Ursache den Gottesdienst alle sehen und erkennen sollten. Denn obwohl Gott die Juden allein zum Eigenthum hatte auserwählet; hat er dennoch immerdar auch etliche aus den Heiden zu sich gezogen. Auf daß sein Volk bald merken und gewöhnen sollte, daß Gott auch der Heiden Gott sei, und sie durchs Evangelium wollte zum Reich Christi bringen.

Denn daß die Königin aus dem Reich Arabia Salomonis Weisheit rühmet, V. 6., ist nicht allein zu verstehen von weltlicher Weisheit, sondern daß er sie aus Gottes Wort, von Gott und dem rechten Gottesdienst hat unterrichtet, um welcher Ursachen willen sie am meisten gen Jerusalem kommen ist, ihn zu sehen und hören, wie man hie sehen kann, 1 Kön. 10, 1., da so geschrieben stehet: Da das Gerücht Salomo, von dem Namen des HErrn, kam für die Königin vom Reich Arabia.

\section{Votum}

Helfe und regiere uns Gott alle, daß wir seine Wahrheit in seinem heiligen Wort auch hören, und in seiner Erkenntniß täglich wachsen und zunehmen, und allzeit den rechten Gottesdienst verrichten, daß sein Name geheiliget und sein Wille geschehe, zu unserm Heil und Seligkeit, um Christi Jesu willen, Amen.

\section{Der andere Theil des 10. Kapitels}

\section{Vorrede}

Daß die Regenten Tugend und Gaben haben, die zu glückseliger Regierung von nöthen sind, das kömmt allein von Gott, dem unerschöpflichen Brunnquell aller Weisheit und Gütigkeit. Solches bezeuget Salomon aus eigener Erfahrung, und sagt an Gottesstatt, Spr. Sal. 8, 14. 15.: Mein ist beide Rath und That, ich habe Verstand und Macht, durch mich regieren die Könige, und die Rathsherren setzen das Recht, durch mich herrschen die Fürsten und alle Regenten auf Erden. Derwegen laßt uns Gott herzlich bitten, er wolle allen christlichen Regenten, als seinen und seines Reichs Amtleuten, zur Verwaltung seines Reichs und Amts, Rath und That geben, auf daß sie allzeit also regieren und herrschen, daß es zu seines heiligen Namens Ehre und zu gemeinem Nutz und Friede diene und gelange, Amen. Darauf höret ferner an den andern Theil des 10. Kapitels.
\quellseite{516}

\section{Summarium}

Sie ist eine Historia, wie Salomo ihm habe einen Königs-Stuhl gebauet, und scheinet fast, als habe er damit, gleich wie durch ein Gemälde, wollen anzeigen die Art und Eigenschaft eines vernünftigen Regenten.

Das Elfenbein, wenn mans recht zubereitet, ist es schön, weiß und glatt, bedeutet derhalben die Unschuld. Denn ein Regent soll erstlich fromm und unbefleckt sein, und für allem sich hüten, daß ärgerlich ist.

Das Gold, damit das Elfenbein überzogen ist, bedeutet die Beständigkeit, daß man sich weder Gunst noch Haß zu Unbilligkeit lasse bewegen, und fest und steif bei dem Rechten halte.

Die sechs Stufen am Stuhl bedeuten, daß ein Regent nicht zu jäh[e] sei, sondern ihm der Weile nehmen soll, auf daß niemand übereilet und er selbst durch eigen Gutdünken nicht verführet werde. Denn zu solchem Amt gehöret, daß man wohl besonnen sei.

Dazu dienet auch, daß der Stuhl hinten rund ist. Denn wo man bedächtig handelt, da vergreift man sich desto weniger.

Die zwo Lehnen, darauf man sich steuret, bedeutet das Recht an ihm selbst.

Die Löwen auf den Lehnen und Staffeln bedeuten das rechte Löwenherz, daß ein Regent muß haben, daß er ohne alle Scheu strafe, wo Strafe vonnöthen ist.

Daß nun Salomo so großen Reichthum gehabt, V. 23., ist ein Zeugniß, wo fromme und gottselige Regenten sind, daß Glück und Heil die Fülle werde da sein. Dagegen findet sich eitel Unglück, wo gottlose und ungeschickte Regenten sind.

\section{Votum}

Gott, der Vater des Lichtes, von dem alle gute Gaben von oben herab kommen, segne und mehre auch heute alle Obrigkeit, an Verstand, Tugenden und Reichthum, und gebe zu ihrer Regierung ferner alle glückselige Wohlfahrt, dem Kirchen- und Haus-Regiment zu Schutz und Erhaltung. Amen.

\chapter{Das 11. Kapitel}

\section{Vorrede}

Gleich wie die Gottesfurcht ist die fürnehmste Tugend, welche in sich begreift die erste Tafel der zehen Gebote Gottes; und wie der 111. Psalm V. 10. saget, der Weisheit Anfang, oder der Anfang und das Ende unserer Bekehrung zu Gott: also ist die Gottesfurcht der Ursprung und die rechte Probe aller Tugenden, die in der andern Tafel der zehen Gebote begriffen werden. Und weil diese Tugend sonderlich ins erste Gebot gehöret, die Gott von unsern Herzen erfordert, so haben die Alten gar fein ein Gleichniß von dem ersten Gebot geben wollen, da sie gesagt: Das erste Gebot mit dieser Tugend der Gottesfurcht ist gleich wie die Schiene an dem Kranze der Braut Christi, auf welche die schönen Blumen und Röslein aller andern Tugenden in den nachfolgenden Geboten müssen aufgebunden werden. Denn gleich wie man seinen Kranz aus Blümlein und Röslein machen kann, wo es an einer Schienen mangelt: also machen auch die Tugenden der andern Tafel keinen Christen, wo nicht kindliche Gottesfurcht vorher geht. Derwegen wir alle in Gottesfurcht uns fleißig die Tage unsers ganzen Lebens üben sollen. Dies unterrichtet auch das 11. Kapitel.

\section{Salomos Weiber, Abgötterei, Kinder und Tod}
\quellseite{517}

\section{Summarium}

Hie ist eine historia, wie Salomo erstlich wider Gottes Befehl und Warnung heidnische Weiber nimmt, V. 1., und hernach ihrem Exempel nachfolget und Abgötterei mit ihnen treibet. Auf daß jedermann lerne Gott fürchten und sich in guter Hut halte, sintemal es so leichtlich geschehen kann, daß auch große Heiligen fallen. Denn wie gehört, ist Salomo ein so hochbegnadeter und in Gottes Wort wohl geübter und erfahrener Mann gewesen, und dennoch fällt er so greuslich
\quellseite{518}
 und gefährlich, daß er seinen Weibern folget und abgöttisch mit ihnen wird.

Zum andern stehet man hie, wie Gott Abgötterei strafet mit Krieg und Zerrüttung des Regiments. Darum sonderlich vorzunehmen ist, daß weltliche Obrigkeit den rechten Gottesdienst handhabe und fördere, und aller Abgötterei mit Ernst wehre.

\section{Summarium}

Hie ist eine Historia, in welcher man klar siehet, wie Gott die Abgötterei also strafet, daß er das herrliche und schöne Königreich seines Volks, das unter David und Salomo so hoch gestiegen war, durch Aufruhr zerrütten und zertheilen läßt. Auf daß weltliche Obrigkeit sonderlich Abgötterei und falschen Gottesdienst lerne fliehen, als das schädliche Gift, das alle Regimente wüst macht, und allen Jammer erregt. Wiederum soll jedermann allen Fleiß dahin wenden, daß der rechte Gottesdienst sammt der reinen Lehre gefördert und gehandhabt werde. Denn wo das geschiehet, so gibt Gott Glück und Heil zum Regiment. Wir sollen aber solche Historia nicht allein zur Lehre, sondern auch zum Trost brauchen. Wehe thut es, daß die Bischöfe und etliche wenige Fürsten wider das Evangelium toben, die armen Pfarrherrn, so Gottes Wort wollen pflanzen, vertreiben und würgen, und an den Unterthanen alle Tyrannei üben. Aber lerne hie, sie können ihr eigen Verderbniß
\quellseite{519}
nicht baß fördern, denn auf diese Weise. Herodes wäre des Herrn Christi, und sein Sohn des heiligen Johannis des Täufers halben wohl König blieben, und in allen Würden gesessen; aber da der Vater die unschuldigen Kindlein und der Sohn den heiligen Johannem läßt würgen, eilen sie nach ihrem Verderben und müssen darum zu Boden gehen.

\section{Votum}

Der starke, allmächtige GOtt laß uns ja den Teufel und unser Fleisch und Blut von seinen Geboten nicht abführen, und regiere uns, daß wir uns für Abgötterei und allen andern Sünden hüten, und steure allen, so muthwillig die Sünden vertheidigen und Gottesfurcht in den Gläubigen verfolgen, Amen.

\chapter{Das 12. Kapitel}

\section{Vorrede}

Die Tyrannei und Abgötterei sind die größten und schädlichsten Sünden der Obrigkeit. Durch Abgötterei wird Gott gelästert und die wahre Religion unterdrückt, zu Schaden der Unterthanen Seelen. Durch Tyrannei muß die Gerechtigkeit Noth leiden, und werden die Unterthanen verderbet. Derwegen laßt uns Gott abermal danken, daß wir gnädige und gütige Obrigkeit, und auch, daß wir von Abgötterei erlöset, Gottes reines Wort und rechten Gottesdienst haben, daß wir ohne Hindernis zur Kirche gehen, mit Lehre, Lernung, Gebet und mit Gebrauch der heiligen Sacramente, Gott loben und preisen können; und daß sich ein jedes in seinem Handel, Acker und Handthierung wohl nähren, und davon durch Gottes Segen das meiste und beste Theil behalten kann. O diese Herrlichkeit und Schätze erhalte uns Gott allzeit, bis an den lieben jüngsten Tag, Amen. Darauf höret davon an das 12. Kapitel.

\section{Zertheilung des Reichs Salomo}
\quellseite{520}

\section{Summarium}

Sie ist eine sehr feine Historia, in welcher man siehet, wie die Tyrannei, zu Aufruhr Ursache gibt, wenn man die Leute zu hart dringen und sonderlich im angehenden Regiment will zu strenge sein, V. 10. Aber junge hitzige Leute, wie Rehabeam, und sein junger Rath denken, sie wollens mit Gewalt hinaus führen, wie sie es haben fürgenommen. Aber es findet sich viel anders.

Sie muß man aber nicht vergessen, daß der Text meldet, daß solcher unzeitiger und hitziger Rath, der zu Aufruhr Ursache gibt, eine Strafe sei der Abgötterei und des falschen Gottesdienstes, welchen Salomo hatte angerichtet.

\section{Summarium}

Sie ist eine sonderliche Historia, wie Jerobeam, nachdem er ins Regiment kommt, besorget, daß Königreich werde ihm nicht bleiben, wo er die zehen Stämme Israel lasse hinauf gen Jerusalem ziehen, und daselbst nach Gottes Befehl opfern. Darum erdenket er aus seinem Herzen einen neuen Gottesdienst; aber eben
\quellseite{521}
damit er sein Reich gedenkt zu erhalten, verderbet er’s, V. 28. Gleich wie heutiges Tages viel sind, die also gedenken: Wenn sie bei dem alten papistischen Wesen bleiben, so wollen sie mancher Gefahr und Unrath entlaufen. Aber damit versündigen sie sich je länger je mehr, bis sie endlich gar zu Boden gehen müssen. Denn wer ein beständig Glück will haben, der denke und halte sich zu Gott und seinem Wort, sonst soll und muß er zu Boden gehen, wenn er gleich aller Kaiser Geld und Gewalt hätte.

\section{Votum}

Der treue Gott behüte die ganze Christenheit gnädiglich für tyrannischem Regiment, und für öffentlicher und heimlicher Abgötterei; und helfe, daß Obrigkeit und Unterthanen ihm, nach seinen Geboten, gehorsamen, zu zeitlicher und ewiger Wohlfahrt, Amen.

\chapter{Das 13. Kapitel}

\section{Vorrede}

Wer da sündiget, der thut übel und unrecht; wer aber gesündiget hat, und sich Gottes Wort strafen läßt, und Buße thut, der thut wohl und recht. Derwegen verharre niemand halsstarrig in seinen Sünden, sondern, da ihm Buße von der Kanzel oder im Beichtstuhl geprediget wird, so erkenne er seine Sünde, lasse die ihm leid sein und bekehre sich zu Gott im wahren Glauben, mit beständigem Vorsatz, daß er wolle fromm sein. Denn wer in Sünden vorsätzlich verharrt, und flapp auf flapp schlägt, das ist, eines nach dem andern sündiget, dem gereicht es endlich zu zeitlichem und ewigem Leibes- und Seelen-Schaden. Denn der Tod ist der Sünden Sold, Röm. 6, 23. Das ist, die Sünde und Bosheit besoldet, belohnet und bezahlet Gott mit allerlei Strafe, mit dem zeitlichen und ewigen Tode oder Verdammniß. Dies soll auch erinnern der erste Theil des 13. Kapitels.

\section[Jerobeams Hand verdorret. Ein Prophet wird vom Löwen]{Jerobeams Hand verdorret. Ein Prophet wird vom Löwen getödtet}

\section{Summarium}

Dies ist eine Historia, wie Jerobeam seiner Abgötterei halben von einem Propheten gestraft wird, V. 2. Da siehest du, daß aller Abgöttischen und ketzer Art ist, daß sie wollen ungestraft sein und Recht haben. Derowhalben befiehlt der König auch, man soll den Propheten greifen, V. 4. Solche Gefahr müssen alle die auf sich nehmen, welchen Gottes Wort zu predigen und Abgötterei und Aergerniß zu strafen befohlen ist.

Neben dem siehet man auch, wie an einem verkehrten Menschen gar nichts hilft. Jerobeam höret den Propheten selbst, und konnte wohl gedenken, daß er von Gott zu ihm gesandt war, erstlich dadurch, da er nach ihm greift, seine Hand ihm so bald dürr wird, V. 4. Darnach auch, da der Prophet für ihn betet, daß ihm die Hand wieder zurecht gebracht wird, V. 4. Daran kehret er sich aber gar nichts, sondern bleibet
\quellseite{522}
bei seinem abgöttischen Wesen wie vor, bis er von der Erde vertilget wird. Wie man heutiges Tages dergleichen Exempel sehr viel erfähret an den Feinden des Evangelii, und täglich mehr erfahren werden. Denn Gott hat nie keines geschonet, es sei denn rechte Buße erfolget.

\section{Votum}

Helfe GOtt, der sein heilig Wort zu lehren und lernen befohlen hat, daß die Lehrer seines Worts die Sünden ernstlich und fleißig strafen, und regiere uns, daß wir uns alle aus den Straf-Predigten recht und wohl bessern, und ihn zeitlich und ewig mit einander loben und preisen, Amen.

\section{Der andere Theil des 13. Kapitels}

\section{Vorrede}

Wer von Gottes Wort abweicht, von dem weicht Gott wieder mit seinem Schutz und Segen; und wer Gottes Wort nicht folget, dem folget dagegen Unglück und Untergang. Das wird täglich gelehret, und davon gibt Gott oft Zeugniß mit wunderbaren und großen Strafen. Dennoch sind ihrer viel, die mehr anderer Leute Rath und ihrem Gutdünken folgen, denn Gott und seinem Wort und Geboten. Derwegen wollen wir unsern Rathgeber allein sein lassen das heilige Wort Gottes, Ps. 119, 24. Und wollen dem gehorsamlich folgen, auf daß wir in Gottes Segen und Schutz unter aller Gefahr geleitet und erhalten werden. Und dazu sollen alle fromme Herzen täglich beten aus dem 119. Psalm V. 4. 29.: HErr, du hast geboten fleißig zu halten deine Befehle; wende von mir den falschen Weg, und gönne mir deine Gesetze. Darauf höret auch an den andern und letzten Theil des 13. Kapitels.
\quellseite{523}

\section{Summarium}

Ist eine wunderbare Historia, wie der Prophet, welcher so eine herrliche Prophezeiung wider Jerobeam und den Altar und hernach ein herrliches Wunderwerk gethan hatte, von einem Löwen zerrissen wird, darum, daß er wider den Befehl Gottes zu Bethel gegessen hatte.

Da werden wir erstlich vermahnet, dem mit allem Ernst zu glauben und fest daran zu halten, was Gott einmal geboten und befohlen hat, und uns nicht davon lassen abwenden und reißen, wenn gleich jemand käme und gäbe für, er wäre ein Prophet, mit dem ein Engel durch des HErrn Wort geredet hätte etc. Denn was hat diesen lieben Mann Gottes anders zu diesem Fall gebracht, denn daß er gemeint hat, weil der andere so herrlich rühmet, er wäre auch ein Prophet, so müßte er Gott gehorsam sein, und ist also durch Gottes Namen betrogen worden.

Darnach hat auch Gott dem Könige Jerobeam und dem ganzen Israel durch solch Werk wollen anzeigen, wie ihnen ihre Abgötterei gerathen soll. Und wie eine große Sünde es sei, da sie ihnen Gottes Wort haben predigen lassen, und sind ihm doch nicht nachkommen, daß sie auch also sollen umkommen und zerrissen werden; denn Gott will, daß man seinem Wort folgen oder alles Unglücks gewarten soll.

\section{Votum}

Erleuchte und regiere uns, GOtt, unser lieber Vater, daß wir durch anderer Leute Schaden klug werden, und ja in all unserm Thun und Lassen unserer Beruf und Stände im Gesetze des HErrn wandeln, und allein GOttes Wort folgen, Amen.

\chapter{Das 14. Kapitel}

\section{Vorrede}

Wie der Götzendiener Hauptmann, der Teufel, Adam auch dahin brachte, daß er mehr von der Kraft eines Apfels hielt, denn von Gott und allen seinen Gaben, die er ihm hatte verliehen, und meinte, er wollte in dem Apfel alle göttliche Weisheit finden, welches Abgötterei war: also hanget uns noch heutiges Tages allen an diese Unart, daß wir von Natur nichts anderes sind, denn lauter abgöttische Menschen und Götzendiener, Röm. 6, 17. Denn was wir selber von uns und unserm Thun halten, das wissen wir wohl; wenn wir es nur wollten bedenken und uns lassen leid sein. Denn daher kömmt alle Abgötterei, daß ein jeder seinem Gutdünken mehr, denn dem klaren Wort Gottes folget, und sich auf die Creaturen mehr, denn auf seinen Schöpfer verläßt. Ach, so laßt uns alle Abgötterei fliehen, und zu Herzen nehmen diese ernste Dräuung und Warnung, als das Moses sagt in seinem fünften Buch im 8. Kapitel V. 19.: Wirst du des HErrn, deines Gottes, vergessen und andern Göttern nachfolgen, und ihnen dienen, und sie anbeten, so bezeuge ich heute über euch, daß ihr umkommen werdet. Und im 14. Kapitel des Buchs der Weisheit saget der heil. Geist V. 27.: Den schändlichen Götzen dienen ist alles Bösen Anfang, Ursache und Ende. Deß ist ein Exempel das 14. Kapitel.

\section[Weissagung wider Jerobeam. Abgötterei, Regirung; Sünde und derselben]{Weissagung wider Jerobeam. Abgötterei, Regirung; Sünde und derselben Strafe}
\quellseite{524}
ᵃ Mimesis, in futuro dicetur, his impletis. Was hat Jerobeam nun gemacht.

\section{Summarium}

Dies ist die dritte Warnung, daß Gott den König Jerobeam seiner Abgötterei halben durch den Propheten Ahia strafet. Aber es ist alle Warnung und Strafe an den Gottlosen verloren, man sage ihnen, was man wolle, so thun sie, was ihnen gefällt.

Und ist hie sonderlich die Predigt des Propheten Ahia zu merken, wie heftig er dräuet, daß Gott die Abgötterei Jerobeams nicht allein an seinem ganzen Hause, sondern auch an dem ganzen Volk Israel strafen wolle, v. 7. 10. Auf daß jedermann sich für Abgötterei hüte, und mit allem Fleiß sich zum rechten Gottesdienst halte.
\quellseite{525}

\section{Summarium}

Ist ein Exempel, wie das Königreich Juda auch Abgötterei treibet, und an andern Herrn, denn zu Jerusalem, opfert. Um solcher Sünde willen erwecket Gott den König aus Egypten, der plündert den Tempel und führet den Schatz daraus mit sich etc. Uns zur Warnung, wenn wir mit Frieden wollen leben, daß wir uns Gottes und seines Worts halten sollen. Wo man aber unrechte Lehre und Abgötterei hat, da soll Krieg und alles Unglück sich mit Haufen finden.

\section{Votum}

Wir danken Gott, dem HErrn, dem Vater aller Gnaden und Barmherzigkeit, daß wir sein heilig Wort und rechten Gottesdienst haben; der behüte uns ferner für aller Abgötterei, und regiere uns, daß wir ihn über alle Dinge fürchten, lieben und vertrauen, bis zu unserm seligen Ende, Amen.

\chapter{Das 15. Kapitel}

\section{Vorrede}

Die zeitlichen Strafen, damit der gerechte Gott Abgötterei und falsche Lehre und Gottesdienst heimsucht und straft, sind schwer; aber am allerschwersten sind die ewigen Strafen. Von denen predigt also die heilige Schrift 1 Cor. 6, 9.: Weder die Hurer, noch Abgöttischen, werden das Reich Gottes ererben. Und in der Offenbarung Johannis im 21. Capitel V. 8. wird gesagt: Der Abgöttischen und aller Lügner Theil wird sein in dem Pfuhl, der mit Feuer und Schwefel brennet. Demegegen, daß nicht zeitliche und ewige Strafen auf uns kommen, so laßt uns herzlich alle Tage beten, daß uns Gott für Abgötterei behüten und sein Wort und rechten Gottesdienst bei uns erhalten wolle. Davon höret auch an das 15. Kapitel.

\section[Regierung zweier Könige in Juda, Abiam und Asa;]{Regierung zweier Könige in Juda, Abiam und Asa; und zweier in Israel, Nadab und Baasa}
\quellseite{526}

a Miplezeth) wer der Miplezeth gewesen, ist ungewiß. Etliche sagen, er sei der Abgott Priapus gewesen.

\section{Summarium}

Dies ist eine Historia, in welcher uns zwei schöne Exempel werden angezeigt. Erstlich, daß der König Abiam nur drey Jahr regiert zu Jerusalem, V. 1.; sein Sohn Asa aber ein und vierzig, V. 10. Dieser Asa hat ein sonderlich fürstlich Werk geübet, daß er wehret dem abgöttischen Wesen und den Abgott Miplezeth ausrottet, V. 13. Denn einem Regenten (soll er anders Glück haben) stehet sonderlich zu, daß er darob sei, daß Abgötterei und falsche Lehre nicht geduldet werde. Denn wo Abgötterei und falsche Lehre ungehindert geht, daselbst schicket Gott zur Strafe Krieg und alles Unglück, wie man hie höret mit Baesa, dem Könige Israel, V. 17.
\quellseite{527}

\section{Summarium}

Ist eine Historia, in welcher uns fürgehalten wird, wie Gott mit den gottlosen Regenten, so Abgötterei und falsche Lehre handhaben, pflegt umzugehen. Jerobeams Sohn, Nadab, regiert nur zwo Jahr über Israel und wird erstochen, V. 25. Baesa, welcher an seine Statt König ward, wird mit Krieg überfallen vom Könige aus Syrien, und endlich sammt seinem ganzen Geschlecht ausgerottet, zur Strafe seiner Abgötterei, und daß er den König Nadab erwürget hat. Also ist dieß Buch durchaus ein rechter Spiegel, darin sich alle Regenten mit allem Fleiß schauen sollen, auf daß sie in Gottesfurcht wandeln und also dem Unrath entfliehen mögen, der über die gottlosen Regenten allweg kommen ist.

\section{Votum}

Der treue GOtt verleihe, daß alle Obrigkeit sein Wort liebe und fördere, und der Abgötterei und falscher Lehre steure und wehre; und regiere uns, daß wir alle Abgötterei fliehen und durch Fürbitz uns und den Unsern nicht Strafen zuziehen, Amen.

\chapter{Das 16. Kapitel}

\section{Vorrede}

Wie Gott das jüdische Königreich wegen der Abgötterei gestraft hat: also wird diese letzte Welt auch ein Schwereres noch ausstehen müssen, wegen der heimlichen und öffentlichen Abgötterei, die bisher getrieben worden und noch im Schwange gehet. Derwegen, daß Gott in gemeinem Trippel und Strafen unser gnädiglich wahrnehme, so laßt uns Buße thun in Reue und Leid über unsere Sünde, in wahrem Glauben und in Gehorsam gegen Gottes Geboten leben, und beten mit dem Könige Zedekia: Das Wort des HErrn ist gut, das du sagest; es sei nur Friede und Treue, weil ich lebe, 2 Kön. 20, 19. Wir wollen aber gleichwohl auch beten, daß Gott an unsern Nachkommen seiner Gnade und Barmherzigkeit nicht vergesse; und wolle ihm allzeit einen heiligen Samen, und ein Häuflein, das ihn erkenne, lobe und preise, und ihm diene in Heiligkeit und Gerechtigkeit, unter ihnen erhalten. Davon höret auch an das 16. Kapitel.

\section{Von vier Königen in Israel: Ela, Simri, Amri und Ahab}
\quellseite{528}

\section{Summarium}

Ist eine Historia, wie Gott das Königreich Israel um der Abgötterei willen mit Aufruhr strafet, V. 9., neben dem, daß die gottlosen Könige einer nach dem andern erstochen wird, um ihrer Sünden willen, daß sie die Abgötterei nicht wehreten, sondern mehreten, V. 10. Zum Exempel aller Obrigkeit, so sie ihr selbst glückselige Regierung und der Unterthanen Wohlfahrt
\quellseite{529}
 begehren, daß sie sich hüten für Abgötterei und Gott und sein Wort für Augen haben, und den rechten Gottesdienst (unangesehen, was andere thun, oder ihnen darüber begegnen mag) mit Ernst fördern.

\section{Votum}

Verleihe GOtt, daß sein heilig Wort unserer lieben Obrigkeit fürleuchte, auf daß ihre Regierung, wie bisher, also ferner, allzeit zur Erhaltung und Förderung seiner Kirchen und den Unterthanen zu Heil und Seligkeit in Christo JEsu, unserm HErrn, diene, Amen.

\chapter{Das 17. Kapitel}

\section{Vorrede}

Wer Gott ehret, den will Gott wieder ehren und ernähren, und wenn die Noth noch so groß wäre. Denn Gott über allen menschlichen Verstand die Seinen erhalten kann, etwa ganz und gar ohne Mittel, etwa durch geringe und kleine Mittel, wie Gottes Wort und die Erfahrung bezeuget. Derwegen sollen wir thun, als die lieben Kinderlein zu thun pflegen: dieselben verrichten mit Fleiß, was sie ihre Eltern heißen, und sorgen nicht, was sie werden zu essen und zu trinken haben, sondern sie befehlen diese Sorge den Eltern und denken, wenn die Zeit des Essens herbei kommen, so werden die Eltern wohl darauf gedacht haben; also sollen wir dem fleißig nachgehen, was uns Gott, unser himmlischer Vater, zu thun befohlen hat, sollen an seinen lieben Sohn glauben, in Gehorsam seiner heiligen Gebote leben, treulich und emsig unsern Beruf treiben, und auch Gott dafür sorgen lassen, wie er uns darinnen ernähren und erhalten werde. Ein jeder diene auch seinem Nächsten in seinem Stande und mit seinem Vermögen; so wird ihm Gott auch dienen in seinem Nothstande und Anliegen, mit Kraft und väterlicher Treue. Dies soll auch erinnern das 17. Kapitel.

\section[Elia in der Theurung von Raben gespeiset; wecket]{Elia in der Theurung von Raben gespeiset; wecket zu Zarpath einen Todten auf}
\quellseite{530}

\section{Summarium}

Hie gehet die Historia an von dem großen Propheten Elia, welchen Gott zu der selbigen Zeit in Israel hat erwecket, da die Abgötterei am größten war. Eben wie bei uns das Evangelium ist an Tag kommen, da das Pabstthum und der falsche Gottesdienst am mächtigsten war. Solches alles thut Gott darum, weil er die Abgötterei endlich strafen will und muß, daß doch die Frommen gewarnet und erhalten werden. Es ist aber hie wieder ein sonderlich Exempel, daß Gott mit einer Theurung länger denn drei Jahr das Land um der Abgötterei willen gestraft habe, V. 1. In solcher Noth gedenkt Gott an die Seinen, wie hie die Raben dem Elia Brod und Fleisch zutragen, V. 6. Und der armen Wittwe durch Eliam ihr Mehlfaß und Oelkrug also gesegnet wird, daß sie die ganze Theurung über genug hatte, V. 14. Denn Gott will den Theil mit der That und Wahrheit behalten, daß er seine Gläubigen wolle ernähren in der Theurung.

\section{Votum}

Wir danken GOtt, unserm lieben Vater, der uns speiset und ernähret; der wolle uns ferner für Theurung und andern wohl verdienten Strafen behüten, und wolle die seinen an Leib und Seele, Gut und Ehre, die seiner Kirchen Gliedmaßen in ihrer Noth sich treulich annehmen, Amen.

\chapter{Das 18. Kapitel}

\section{Vorrede}

Der erste Theil des 18. Kapitels des ersten Buchs der Könige hält uns für mancherlei Trost und Lehre. Der eine Trost ist: Wenn Gott die Gottlosen und Bösen mit Strafen im Zorn und Grimm heimsucht, so erhält er oft väterlich und wunderlich die Gläubigen und Frommen. Der andere Trost ist, davon auch Christus lehret im 11. Kapitel Johannis, des heiligen Evangelisten, da er sagt V. 9.: Wer des Tages wandelt, der stößet sich nicht, das ist, wer seines Berufs abwartet, im Glauben und mit Gebet zu Gott, der gefällt Gott wohl und dem hilft er in aller Gefahr. Und da er schon etwas Noth und Widerwärtigkeit leidet, dennoch überwindet er endlich und seine Lehre, Rath und Thaten sind dem Nächsten nütze und heilsam. Die eine Lehre aber ist: Wer von Gott gestraft wird, der lege die Schuld nicht auf andere Leute, sondern erkenne, daß er solche Strafe mit seinen Sünden selber verdienet habe, und thue wahre Buße, so wird ihm Gott auch aus Barmherzigkeit helfen. Die andere Lehre ist: Ein jeder Christ soll Gott herzlich anrufen, daß er uns mit Hunger des Leibes oder nach Brod, oder auch nicht mit Hunger nach seinem heiligen Wort zu hören, nicht strafe, wie er also gestraft hat den König Ahab, und den größten Theil des Volkes im israelitischen Königreiche, sondern daß er dem Leibe geben wolle das tägliche Brod, und der Seelen auch ihre Speise, nämlich das Wort, das aus seinem Munde gehet, und das eine Kraft ist, die da selig macht alle, die daran glauben.

\section{Der Prophet Elia schlachtet die BaalsPfaffen}
\quellseite{531}

\section{Summarium}

Dies ist erstlich ein fein Exempel, wie ein frommer Mann Obadja am Hof des gottlosen Königs Ahab ist, V. 3., durch welchen hundert Propheten erhalten und heimlich in der großen Theurung ernähret werden, V. 4., auf daß man sehe, wie Gott durch mancherlei Mittel zur Zeit der Noth die Seinen erhalten kann.

Zum andern ist sonderlich zu merken, wie die Gottlosen alles Unglück, das ihnen ihrer Sünde halben widerfähret, auf die Frommen legen. Denn Ahab gibt hie dem Propheten Elia die Schuld, daß Theurung und allerlei Unrath im Lande ist, V. 18. Eben wie zu unserer Zeit die blinden verstockten Papisten das liebe heilige Evangelium schänden und lästern, in dem, daß sie ihm Schuld geben, als sollte der Bauern Aufruhr und ander Unglück, Ketzerei, Krieg, Theurung daraus entstanden sein etc. Aber Elias spricht: Es sei solches alles der Sünden Schuld, daß Ahab und das Volk Gottes Gebot verlassen und dem Baal haben nachgewandelt.

\section{Votum}

Der treue Gott sei uns gnädig und barmherzig, und helfe, daß ein jeder Mensch seine Sünde erkenne, die Gott abbüße, davon ablasse, und nach Gottes Geboten in allem Gehorsam lebe und wandele, Amen.

\section{Der andere Theil des 18. Kapitels}

\section{Vorrede}

Die Prediger, die recht lehren und wohl leben, die springen nicht allerdings aus einem Steine oder Busche, wenn man dran schläget, wie die gottlose Welt spottet; sondern sie sind gewiß eine Gabe Gottes. Darum denn unser lieber Gott höchlich zu bitten ist, wie der Herr Christus selber lehret Luc. 10, 2., da er sagt: Die Ernte ist groß, der Arbeiter ist wenig; bittet den Herrn der Ernte, daß er Arbeiter aussende in seine Ernte. Der Herr dieser Ernte ist der Herr Christus selbst, der seinen Samen des Worts läßt ausstreuen hie auf Erden, ihm eine reiche Ernte unter dem menschlichen Geschlechte anzurichten; die Arbeiter sind alle treue Lehrer und Seelsorger, welchen, da sie den Samen göttlichen Worts unter den halsstarrigen und widerstrebenden Leuten sollen ausstreuen, es wohl so sauer wird, als irgend einem Schnitter oder Ernten-Knechte. Daß der Herr Christus solche Arbeiter aussende in seine Ernte, soll man bitten, das ist, daß er sie gebe und insinuire mit seinem heil. Geist von oben herab, und lege ihnen selbst das Wort in den Mund, damit sie Frucht und Nutzen schaffen. Helfe nun Gott, daß wir diese Erinnerung wohl in Acht nehmen und recht gebrauchen. Amen. Darauf höret ferner an den andern Theil des 18. Kapitels.
\quellseite{532}

a Hinkten die falschen Heiligen, wenn sie andächtig sein wollen, werfen sie den Kopf zu beiden Seiten, wie eine Gans gehet, daß es scheinet, wie ein hinkender zu beiden Seiten hin und her wackelt. Daher auch das Paßah kommen, 2 Mos. 12, 12. Da GOtt in Egypten durchging, und hinket, das ist, er schlug zu beiden Seiten, hie und da, wie ein trunkener Mann gehet. Sic supra: Ut quid claudicatis in duas partes? Significat impios esse vagos, et duplices animo, nihil certi habere, Eph. 4, 14. Inde etiam saltare dicuntur tales, quia motu capitum similes sunt salientibus et vere capitibus saltant. Das ist: so oben: Was hinket ihr auf beiden Seiten? zeigt an, die Gottlosen seien wankend und von doppeltem Herzen, haben nichts Gewisses, Eph. 4, 14. Daher sagt man auch von solchen, daß sie tanzen, weil sie mit der Bewegung der Häupter ähnlich sind dem Tanzenden und in der That mit den Häuptern tanzen.
\quellseite{533}

\section{Summarium}

Ist eine schöne Historia, wie Gott dem frommen Propheten Elia beistehet, und gibt ihm ein öffentlich Zeugniß seiner Lehre und Berufes, dagegen aber die Propheten Baals mit Schaden bestehen und werden um der Abgötterei willen alle erwürget. Auf daß man lerne, sich mit Ernst zum rechten Gott und seinem Wort halten, will man anders der leiblichen Strafe und dem ewigen Zorn Gottes und Verdammniß entlaufen.

Daß aber Elias die falschen Propheten selbst erwürget, thut er darum, daß es Gott mit einem öffentlichen Zeichen beweiset hatte, daß sie Verführer und Gotteslästerer wären, die dem Gesetz nach mit dem Tode mußten gestraft werden. Solches, da es der König nicht wollte thun, that es Elias, nicht aus eigenem Fürnehmen, sondern aus einem besondern prophetischen Geist und Muth, den er von Gott hatte. Im Evangelio vergleicht Christus Johannem den Täufer mit dem Elia, nicht daß er auch das Schwert geführet habe, wie Elias, sondern daß er so mit einem hohen Geist und trefflichem Ernst sein Amt geführet, die Hohenpriester und Phar isäer ohne Scheu gestrafet und die Leute zum rechten Gottesdienst gebracht habe.

\section{Votum}

Der treue GOtt regiere uns alle, daß wir von seinem Wort nicht weichen, sondern dem allzeit anhangen und gehorsamlich folgen; gebe auch allzeit christliche Lehrer und Seegen zu ihrem Amt, und steure allen falschen Propheten und Predigern, Amen.

\chapter{Das 19. Kapitel}

\section{Vorrede}

Fromme Zuhörer sind ihren Predigern verbunden, daß sie die hören sollen. Nicht allein aber, wenn die Prediger reden, was den Zuhörern wohlgefällt (wie denn St. Paulus geweissagt hat, daß am Ende der Welt der Zuhörer gar viel sein und eine Zeit sein wird, da sie die heilsame Lehre nicht leiden werden, 2 Tim. 4, 3.); sondern auch, wenn die Prediger wider die Zuhörer reden, sie dadurch ihrer Missethat zu erinnern und zu bekehren: also sollen die Zuhörer ihre Prediger hören, nicht allein, wenn sie lauter Evangelium und Gnade lehren, sondern auch, wenn sie Mosen und das Gesetz predigen. Das drücke ein jedes in sein Herz zum ewigen Gedächtniß. Denn wehe den Zuhörern, welche die Wahrheit von ihren Seelsorgern trennet und scheidet; wie die Galater die liebe Wahrheit, die ihnen zum Besten gepredigt ward, also von St. Paulo trennet, daß sie ihm todfeind worden, Gal. 4, 16. Dies soll auch erinnern das 19. Kapitel.

\section{Elias Flucht, Fasten, Trost und Nachfolger}
\quellseite{534}

\section{Summarium}

Ist eine Historia, in welcher man siehet, wie allweg die Propheten um des Worts und rechten Gottesdienstes willen sich haben leiden müssen; und hilft sie nichts, obgleich Gott mit öffentlichen Wunderzeichen vom Himmel bezeuget hat, daß ihre Lehre und Wesen recht sei. Solches macht die Heiligen oft kleinmüthig; wie Elias hie vor großer Angst und Traurigkeit begehrt zu sterben, V. 4. Aber Gott tröstet sie und hilft ihnen aus aller Noth.

Zum andern ist ein fein Exempel hie, daß Gott die sieben tausend, so Baal nicht gedienet haben, behüten will, V. 19., da er den andern, so dem Baal dieneten, dräuet, sie sollen durchs Schwert Hasael, Jehu und Elisa fallen, V. 17. Denn Gott will die Abgötterei nimmermehr ungestraft lassen.

\section{Votum}

Der barmherzige, gütige GOtt tröste, stärke und erhalte alle treue Prediger und Lehrer in ihrem schweren und gefährlichen Amt; schütze und segne auch alle fromme Zuhörer, und stürze alle Verächter und Verfolger seines heiligen Namens, Wortes und gläubiger Bekenner, Amen.

\chapter{Das 20. Kapitel}

\section{Vorrede}

Zwei nötiger Lehren sollen wir uns heute erinnern. Die erste ist: Wenn Gott die Menschen mit einer kleinen Strafe heimsucht, so folget nicht lange darauf eine große, wenn man nicht wahre Buße thut, und sich zu Gott ernstlich bekehret. Dies bestätiget die Weisheit Gottes in Sprüchen Salomonis, im 1. Kapitel, V. 24., da sie saget: Ich recke meine Hand aus, und niemand achtet darauf, und lasset fahren allen meinen Rath, und wollet meiner Strafe nicht; so will ich auch lachen in eurem Unfall und euer spotten, wenn über euch kommt, das ihr fürchtet. Wenn über euch kommt, wie ein Sturm, das ihr fürchtet und euer Unfall als ein Wetter, wenn über euch Angst und Noth kommt. Die andere Lehre ist: Wer alles haben will, der behält endlich gar nichts. Davon zeuget Sirach im 11. Kapitel V. 10. und spricht: Mein Kind, stecke dich nicht in mancherlei Händel; denn wo du dir mancherlei fürnimmst, wirst du nicht viel daran gewinnen. Wenn du gleich fast darnach
\quellseite{535}
ringet, so erlangest du es doch nicht, und wenn du gleich hie und da fliehest, so kommst du doch nicht heraus. Demgegen laßt uns dieser Lehren gebrauchen zur Buße, daß Gott unser mit seinen Strafen verschone, und zur Warnung, daß ein jeder also das Seine mehre, damit es ihm Gott auch bewahre und erhalte. Davon höret an den ersten Theil des 20. Kapitels.

\section[Zweyfacher Krieg und Sieg Ahabs wider Benhadab, den]{Zweyfacher Krieg und Sieg Ahabs wider Benhadab, den Syrer}

a Nicht rühmen) das ist: Er spreche nicht hui, ehe er über den Berg komme; denn wer genommen hat, der lege den Harnisch ab, und mag sich rühmen; wer ihn aber anlegt, hat drum noch nicht gewonnen.

\section{Summarium}

Sie ist eine Historia, wie Gott immerdar mit der Strafe innen hält, dieselbige nur von ferne weiset und wartet, ob man sich nicht bessern wolle. Also lässet er die Syrer ins Land fallen, aber doch nicht wider Israel siegen, v. 1. Zum sonderlichen Exempel den Tyrannen, wenn sie mehr, denn billig ist, wollen haben, daß sie lernen, daß Gott solche Unbilligkeit strafen wolle. Denn
\quellseite{536}
Ahab wollte den König zu Syrien zu einem Oberherrn erkennen; da aber der König zu Syrien sich an solchem nicht wollte genügen lassen, sondern alles haben, ward solche Tyrannei gestrafet.

\section{Votum}

Gott, unser lieber Vater, regiere uns, daß wir, seine Kinder, durch seine kleinen Ruthen und Heimsuchung fromm werden, auf daß uns nicht große verderbliche Straßen folgen; und helfe, daß ihm ein jeder also an dem Seinen genügen lasse, damit ihn Gott ferner an Gut und Ehre segne, aus Gnaden, Amen.

\section{Der andere Theil des 20. Kapitels}

\section{Vorrede}

Gottes-Lästerung und das Vertrauen auf Menschen-Hülfe sind zwei schändliche und schädliche Laster. Ueber die Gottes-Lästerung klaget der Prophet Jesaias im 1. Kapitel V. 4. und sagt: O weh des sündigen Volkes, des Volkes von großer Missethat, des boshaftigen Samens der schädlichen Kinder, die den Herrn verlassen, den Heiligen in Israel lästern, weichen zurück. Das Vertrauen auf Menschen-Hülfe strafet der Prophet Jeremias im 17. Kapitel V. 5. also: So spricht der Herr: Verflucht ist der Mann, der sich auf Menschen verläßt, und hält Fleisch für seinen Arm, und mit seinem Herzen vom Herrn weicht. Dagegen wollen wir Gottes Allmacht und in der Noth seinen Beistand loben und preisen mit Mund, Herz und gottseligem Wandel, und wollen in allem betrübten Zustande allein auf Gott trauen und bauen, Hülfe bei ihm suchen und auf seine Hülfe allzeit beständig warten. Denn es ist gut, sagt der königliche Prophet David Ps. 118, 8., auf den Herrn vertrauen, und nicht sich verlassen auf Menschen. Dagegen bringet Schaden, auf Fürsten sich verlassen, die Menschen sind und die ja nicht können helfen; wie abermals David spricht Ps. 146, 3. Dies lehret auch der andere Theil des 20. Kapitels.
\quellseite{537}

\section{Summarium}

Dies ist erstlich eine feine Historia, wie die Syrer geschlagen worden, nicht darum, als wäre ihr Gegentheil so fromm, sondern darum, daß die Syrer Gott gelästert und gesagt hatten, der Israeliten Gott würde ihnen nicht helfen können, wenn sie auf der Ebene eine Schlacht hätten. Auf daß man sehe, daß Gotteslästerung nicht ungestraft bleibe.

Zum andern ist eine Historia hie, in welcher man siehet, daß Ahab sich versündiget, daß er des Sieges, welchen ihm Gott gegeben hatte, nicht recht brauchet, V. 34. Er sollte den gottlosen König aus Syrien getödtet haben; so tröstet er sich seiner Macht und Hülfe, und macht einen Bund mit ihm, und vergisset also, daß er sich allein auf Gottes Hülfe sollte verlassen. Darum muß solcher Sünde sein ganzes Geschlecht hernach entgelten, V. 42. Es ist aber solches eine Fürsten- und Regenten-Sünde; die lassen sich bedünken, wo man mit Macht gefasset sei, da habe es nicht Noth. Aber Gott lehret sie anders, daß es nicht an der Macht, sondern an seiner Gnade und Segen liegt.

\section{Votum}

Der treue Gott behüte uns, daß wir ihn und seine Werke ja nicht lästern; und helfe, daß wir in Nöthen allein auf ihn sehen all unsere Zuversicht, und ihn loben und preisen, zeitlich und ewig, wegen seiner väterlichen Hülfe und treuen Errettung, Amen.

\chapter{Das 21. Kapitel}

\section{Vorrede}

Keinem armen Sünder, obgleich seine Sünden blutroth wären, und er des Tages, nicht allein sieben, sondern siebenzig mal sieben mal, gesündiget hätte, wird die Gnaden Thür verschlossen, sondern stehet ihm allezeit offen, wofern er anders den rechten Weg der Buße dazu trifft, sich demüthiget und zu Gott bekehret. Das bezeugen viele trostreiche Sprüche und gnadenreiche Exempel. Ein fürnehmer Spruch ist, daß von Gott Micha, der Prophet, im 7. Capitel V. 18. sagt: Wo ist solch ein Gott, wie du bist? Der die Sünde vergibt und erlässet die Missethat den Uebrigen seines Erbtheils, der seinen Zorn nicht ewiglich behält, denn er ist barmherzig. Er wird sich unser wieder erbarmen, unsere Missethat dämpfen und alle unsere Sünde in die Tiefe des Meers werfen. Ein schön Exempel aber, daß sich Gottes Barmherzigkeit erstrecke über alle Wolken, und erfülle Himmel und Erden, ist der bußfertige Sünder, der König Ahab. Davon das 21. Capitel.

\section{Geduld Tyranni Ahabs und Jesabels wider Naboth}
\quellseite{538}

a Wenn du) Tu nunc facis regnum super Israel id est: Bist du es doch, der es macht, daß Israel ein Königreich ist, ohne dich wäre es ein Nichts.

b Gesegnet) das ist: Er hat gelästert; denn die Schrift das grausame Laster, das Gott lästern heißt, an diesem Ort aufs züchtigste nennet, wie auch Hiob 2, 9. stehet. Gott lästern aber hatte Mose bei dem Tode verboten. So war das Königreich Israel von Gott eingesetzt, daß den König lästern auch Gott betraf.
\quellseite{539}

c Sack) das ist, betrübte, geringe Kleider, als wenn man Leid trägt, oder arbeitet.

\section{Summarium}

Es ist erstlich eine Historia, wie Ahab sammt seinem Weibe, Isabel, Naboth, den Gerechten, ohne alle Ursache, wider Gott und Recht hat lassen erwürgen, darum, daß er dem Befehl Gottes nach bei seinem Erbe bleiben und dasselbige dem Könige nicht verkaufen wollte, V. 3. Denn obgleich weltliche Obrigkeit Macht hat über der Unterthanen Leib und Gut; so soll es dennoch seine Maße haben, daß ein Unterschied bleibe zwischen einem ordentlichen Regenten und Tyrannen. Ein Tyrann nimmt, was ihm gefällt, und fraget nicht, ob ers Fug oder Recht habe. Ein Regent aber bleibt bei seinem ordentlichen Einkommen, und wo es die Noth nicht sonderlich erfordert, legt er keine neue Beschwerden auf. Solche Gleichmäßigkeit hält Ahab mit dem Naboth nicht, sondern führet tyrannisch. Ob nun wohl der fromme Naboth darum zu Boden gehet, so muß doch Ahab auch darum gestraft werden.

Zum andern ist ein sonderlich Exempel hie, daß Ahab, nachdem der Prophet Elias seiner Sünden halben ihn strafet, V. 19., Buße thut und Gott um solchen Demuth willen die Strafe bei Ahabs Leben auffhält, V. 27. Solches Exempel dienet uns dazu, daß wir im Unglück und wenn Gottes Strafe auf uns liegt, nicht sollen verzagen, sondern uns bekehren, vom Bösen ablassen und um Gnade bitten, so will Gott, wie mit Ahab, sich gnädig finden lassen, der doch ein sehr böser Mensch gewesen war.

\section{Votum}

Der barmherzige, gütige GOtt sei uns auch gnädig, vergebe uns unsere großen vielfältigen Sünden, und lindere die wohlverdienten Strafen; er wolle uns auch regieren, daß wir unserm Nächsten an seinem Leib und Gut keinen Schaden noch Leid nimmer thun, Amen.

\chapter{Das 22. Kapitel}

\section{Vorrede}

Daß die Gottlosen für den göttlichen Strafen sich nicht fürchten, ja auch dieselben nicht glauben, wenn sie ihnen angekündiget werden; diese Blindheit kommt vom Teufel, welcher die Gottlosen nicht allein so sicher macht, sondern sie auch verhindert an der Buße, daß sie Gottes Strafen nicht entrinnen mögen. Denn er weiß wohl, daß Gott durch den Propheten Esaías im 45. Capitel V. 22. saget: Wendet euch zu mir, so werdet ihr selig aller Welt Ende, denn ich bin Gott und keiner mehr. Und wie wohl sagen, daß alle, die dem HErrn widerstrecken, müssen zu Schanden werden. Verwegen wollen wir dem Teufel nicht Raum geben, wollen für Sündigkeit uns hüten, nicht in Sünden verharren und verstarren, und wollen Gottes zeitliche Strafen gebrauchen zur Bekehrung und gottseligem Wandel, daß wir nicht härter allhier und dort ewig gestraft werden wegen unserer Sünden. Darauf höret ferner an den ersten Theil des 22. Capitels.

\section{Ahabs Untergang. Josaphats und Ahabs Regierung}
\quellseite{540}

\section{Summarium}

Dies ist eine wunderbare Historia, wie Gott verhänget, daß der König Ahab, so ohne das Gottes und seines Wortes nicht viel achtet, durch falsche Propheten verführet wird und endlich umkommen, wie ihm zuvor durch einen Propheten geweissaget, v. 6. Sein Bruder, der König zu Syrien, mit dem Ahab einen Bund machte und sich seiner Macht tröstete, gibt ihm seinen rechten Lohn, darum, daß er ihn hatte los gelassen. Hie sieht man, wie die falschen Propheten mit ihren Lügen toben, und die Frommen und rechtschaffenen Propheten sich von ihnen leiden müssen. Aber es währet beides nur eine Zeitlang.

\section{Votum}

Gott steure in seinem Zorn allen falschen Lehrern, und gebe den rechtschaffenen treuen Predigern Geduld unter dem Kreuz und Verfolgung, und helfe, daß sie ihr Amt verrichten zu seinen Ehren und ihren Zuhörern zu Nutz und Seligkeit, Amen.

\section{Der andere Theil des 22. Kapitels}

\section{Vorrede}

Die ganze heil. Schrift neben allen bewährten Historien bezeuget, daß allezeit greuliche Sicherheit bei denjenigen in vollem Schwange gegangen ist, welche Gott der HErr mit seinem heiligen göttlichen Wort reichlich besucht und begabet hat; sonderlich aber ist dieses an den Zeiten des Volkes Gottes, der alten Israeliten, zu sehen, in was großer Sicherheit sie, da sie
\quellseite{541}
Gottes Wort überflüssig gehabt, gelebt haben. Daher rücket Moses im 5. Buch im 32. Kapitel V. 15. seinen Juden dieß Laster heftig auf: Da Israel (spricht er) feist und satt ward, wurde er geil. Er ist dick und stark worden, und hat den Gott fahren lassen, der ihn gemacht hat; er hat den Fels seines Heils gering geachtet. Also wenn die Sicherheit unter den Leuten am größten ist, so ist die Strafe Gottes am nächsten mit dem Garaus, das ist, mit dem letzten Untergang und Verderben; als deß ein Exempel ist der König Ahab mit seinem Untergang, und lehret also, daß ein jämmerlich Ende nehmen alle Gottlose und Abgöttische, welche Gott und sein Wort samt treuen Predigern verachten und verfolgen. Derwegen schaffet, daß ihr selig werdet mit Furcht und Zittern. Denn Gott ist, der in euch wirket, beide das Wollen und das Vollbringen. Und dies geschehe ja in seiner Furcht, und nach seinem Wohlgefallen. Davon höret auch an den andern und letzten Theil des 22. und letzten Kapitels.

a. Da das Schwert anhänget von der Achsel über her, bis auf die Hüfte.

b Huren) da die Huren waschen. Das ist, Canponas lotrices.
\quellseite{542}

\section{Summarium}

Ist erstlich eine Historia, wie der fromme Prophet Micha, darum, daß er die Wahrheit sagt, vom gottlosen König Ahab gefangen wird, V. 27. Zum andern, daß Ahab im Kriege umkommt, V. 37. Zum dritten ist eine sonderliche feine Historia, daß Josaphat die Abgötterei, so noch im Lande war überblieben, vollends ausrottet, V. 43. Denn dies stehet sonderlich christlicher Obrigkeit zu, daß sie die Abgötterei nicht dulden, sondern Gott zu Ehren und den Leuten und gemeinem Nutz zu gut, dieselbige mit Ernst abschaffen und strafen sollen.

\section{Votum}

Gott tröste und erhalte alle treuen Lehrer und Seelsorger unter ihrem Kreuze, strafe die Verfolger seiner Kirche und Christenheit; und behüte uns ferner für aller Abgötterei und andern Sünden wider Gott und wider die Menschen, Amen.

\part{Das zweite Buch von den Königen}

\chapter{Das 1. Kapitel}

\section{Vorrede}

Der königliche Prophet David vergleichet Gottes Wort einem Stecken und Stab, Ps. 23, V. 4. Nun dienet ein Stecken oder Stab einem Wandersmann zu vielem. Erstlich, so er gefallen, daß er dran aufstehe. Zum andern, daß er sich dran halte, damit er nicht auf schlüpfrigen Wegen strauchele. Zum dritten, daß er sich damit der bösen Hunde, die ihn anfallen wollen, erwehre. Zum vierten, daß er sich drauf stemme, und dran ruhe, sonderlich da er schwer zu tragen hat: Also dienet uns auch Gottes Wort, erstlich, so wir allhier, auf den schlüpfrigen Wegen gefallen, daß wir dran aufstehen von Sünden, und uns, durch den Glauben an Christum, wieder aufrichten. Zum andern, daß wir uns dran halten, damit wir nicht fallen in Sicherheit, gottlos Wesen und Verzweiflung. Zum dritten, daß wir damit Sünde, Tod, Teufel, Hölle, als den bösen Hunden, einen Streich bieten können. Zum vierten können und sollen wir Gottes Wort unter der schweren Last des Kreuzes gebrauchen und darauf ruhen und uns damit trösten. So helfe Gott, daß uns sein Wort auch ein Stecken und Stab sei, im Leben und Sterben, Amen. Darauf höret nun an das 1. Kapitel.

\section{Ahasias Krankheit. Elias Wunderwerke}

\section{2 Kön. Kap. 1. v. 1}
\quellseite{543}

\section{Summarium}

Ist erstlich eine Historia, wie der gottlose König Ahazia dem Exempel seines Vaters Ahab folget, V. 53. Weil er nun so gottlos war, daß er bei der Heiden Göttern Rath und Hülfe suchet, V. 2., und nicht bei dem Gott Israel, strafet ihn Gott, daß er in seinem Hause fällt und des Falls hernach stirbt. Auf daß wir uns lernen für falscher Lehre, unrechtem Gottesdienst und Zauberei hüten, und desto fleißiger uns zum Wort Gottes und rechtem Gottesdienst halten.

Denn weil die Abgötterei verursachet, daß Gott dadurch gelästert wird, kann sie Gott in die Länge nicht ungestraft lassen. Sonderlich aber an den Regenten, so dazu helfen, welche doch Amts halben mit allem Ernst der Abgötterei wehren sollen.

\section{Summarium}

In eine sonderliche Historia, daß der Prophet Elias heißet Feuer vom Himmel fallen über die zween Hauptmänner des Königs Ahazia, der ein jeglicher fünfzig Mann bei sich hatte. Eben wie der Apostel St. Petrus, aus göttlicher Kraft, allein mit einem Worte Ananiam und sein Weib Saphiram um der öffentlichen Lügen willen tödtet. Diese zween Hauptmänner kommen mit einem Befehl vom abgöttischen Könige, und ist gewiß, daß sie auch abgöttisch gewesen sind. Weil sie aber Eliam einen Mann Gottes heißen, V. 10., und doch den Gott, welchen Elias ehrete und von ihm predigte, nicht für den rechten Gott hielten — sie waren Baals Diener, wie ihr König — will Elias an ihnen mit solchem Wunderwerk, daß das Feuer vom Himmel fällt, beweisen, daß seine Lehre und Gottesdienst sei recht und ihrer Propheten Lehre und Gottesdienst sei unrecht und falsch. Daß also dies Wunderwerk nicht dahin gehet, daß Elias sich räche, sondern daß Gottes Name geehret und der rechte Gott erkennet werde.

\section{Votum}

Der Sohn Gottes, unser HErr JEsus Christus, regiere uns, daß wir allezeit falsche Lehre, unrechten Gottesdienst und Zauberei fliehen, und dagegen, daß wir seinem Wort folgen, ihm in Heiligkeit und Gerechtigkeit dienen, und daß wir ja mit den Gottlosen nicht umkommen und verderben, Amen.
\quellseite{544}

\chapter{Das 2. Kapitel}

\section{Vorrede}

In den Geschichten der heiligen Apostel im 14. Kapitel V. 22. sagt St. Paulus: Durch viel Trübsal müssen wir ins Reich Gottes gehen. So sagt man: Es sei ein gut Kraut, das heißt, achte es nicht, das heilet allerlei Schmerzen, wenn mans drauf leget und wohl darauf erwärmen läßt. Das rechte, achte es nicht, sollen wir aus dem himmlischen Würzgarten des göttlichen Worts nehmen, sollen es legen, in wahrem Glauben, auf unsere Noth, und den heiligen Geist bei uns wohl erwärmen lassen, alsdann wir in Christo getrost das mannigfaltige Unglück nicht achten, sondern werden uns Gottes Hülfe und des ewigen Lebens mit Freuden unter allem Kreuz trösten können. Ferner in der 1. Epistel an die Thess. 4, 8. spricht abermals St. Paulus: Wer verachtet (vernehmt die Prediger göttlichen Worts), der verachtet nicht Menschen, sondern Gott, der seinen heiligen Geist gegeben hat in euch. Ja der Herr Christus selber sagt auch Luc. 10, 16.: Wer euch (verstehet die Lehrer des Worts Gottes) verachtet, der verachtet mich; wer aber mich verachtet, der verachtet den, der mich gesandt hat. Davon erinnert auch das 2. Kapitel.

\section{Elias Himmelfahrt. Elisa's Wunderthaten}

a Von deinen Häupten) zum Häupten sein, heißt Meister und Lehrer sein; zu Füßen sein, heißt Schüler und Unterthan sein. Denn wenn der Lehrer lehret, sitzt er höher, denn die Schüler, daß er sie zu Füßen, und sie ihn zu Häupten haben. Also sagt St. Paulus Apost. 22, 3.: Er habe zu den Füßen Gamaliel das Gesetz gelernt. Und ist fast eine gemeine Weise der Schrift, also zu reden auf Hebräisch.

b Nicht wollte Elisa zwiefältigen Geist Elia haben, so es doch ein Geist ist, 1 Cor. 12, 4. in allem Heiligen, sondern einen zwiefältigen Mund desselbigen Geistes, daß er stärker und mehr predigen könnte, denn Elia, als er auch that.
\quellseite{545}

\section{Summarium}

Dies ist endlich eine Historia, wie Gott den Propheten Elia in einem Wetter von diesem Leben wegnimmt zum ewigen Leben, V. 2. 11. Zum Trost nicht allein jenen, so mit Elia recht geglaubt und den Baal nicht angebetet hatten, sondern auch uns allen, so da warten auf die selige Hoffnung und herrliche Erscheinung des großen Gottes, und unsers Heilandes Jesu Christi. Daß wir gewiß wissen und hoffen sollen, daß uns ein ewig Reich, nachdem wir aus diesem Jammerthal fahren, in jenem seligen Leben bereit sei, darin wir in ewiger Seligkeit mit unaussprechlicher und herrlicher Freude leben sollen.

Zum andern hebt sich hie an die Historia von dem Propheten Elisa, V. 12. Denn weil es mit dem Königreich Israel an dem war, daß es sollte um der Abgötterei willen gestraft werden, läßt ihnen Gott zuvor sein Wort reichlich dessen halben verkündigen. Wie aber eine große Sünde es sei, wo man Gottes Wort und desselben Verkündiger verachtet und unehret, siehet man hie an den Kindern, welche die wilden Bären zerreißen, darum, daß sie des alten Propheten spotten, V. 23., und zur Warnung, daß wir die Kirchendiener, um ihres Amts willen, lernen in Ehren halten.

\section{Votum}

Helfe der treue Gott, daß wir uns als gehorsame und dankbare Zuhörer seines Worts gegen ihn und unsere Seelsorger und Lehrer halten; und verleihe uns, daß wir alles Kreuz in Geduld überwinden, und nach diesem Leben auch das ewige Leben erlangen, um Jesu Christi, unsers Herrn willen, Amen.

\chapter{Das 3. Kapitel}

\section{Vorrede}

Gott sagt durch den Propheten Zephania, 1. Capitel V. 17.: Ich will den Leuten bange machen, darum, daß sie wider den Herrn gesündiget haben. Dieses ist erfüllet am König Joram. Dagegen spricht Sirach, 1. Kap. V. 19.: Wer den Herrn fürchtet, wenn er Trost bedarf, so wird er gesegnet. Das ist erfüllet am Könige Josaphat. Ferner im Ps. 15, 4. stehet: Ein Gläubiger achtet die Gottlosen nichts, sondern ehret die Gottesfürchtigen. Des hat sich gehalten der Prophet Elisa. Und St. Paulus Röm. 13, 2. sagt: Die der Obrigkeit widerstehen, werden über sich ein Urtheil empfahen. Das haben an sich die Moabiter erfahren. Er erzählet Sprüche. Exempel lauten aus dem 3. Kapitel.

\section{Jorams Regierung. Krieg und Sieg wider die Moabiter}
\quellseite{546}

\section{Summarium}

Es ist erstlich eine feine Historia, in welcher man siehet, was für ein Unterschied sei zwischen den Frommen und Gottlosen, zur Zeit der Noth. Die Gottlosen verzagen bald und können sich Gottes Worts nicht trösten, wie die Joram thut, V. 10., denn sie haben es verachtet und nicht gemerket. Aber die Frommen, wie Josaphat, haben Gottes Wort in Ehren, merken und behalten es, und wenn die Noth hergehet, halten sie sich daran und setzen ihr Vertrauen darauf, und finden derhalben reichen Trost darin, V. 12. Solches genießen hernach die Gottlosen auch.

Zum andern ists auch ein schön Exempel mit dem Propheten, daß er dem gottlosen König unter Augen sagt: Er wollte ihn nicht ansehen, wenn er nicht Josaphat ansähe. Warum er nicht seine Baals-Propheten um Rath fraget? V. 13. Solchen Ernst ist ein Prediger schuldig wider gottlose Obrigkeit zu brauchen, sintemal doch kein Ehren noch Vermaynen sonst an ihnen helfen will.
\quellseite{547}

a Vel quod ipsi Israel offensi, ista abominatione timerunt, ne Deus in ipsos quoque irasceretur. D. i. weil die Israeliten, selbst über solchen Greuel erzürnt, fürchteten, Gott möchte auch gegen sie selbst zornig werden.

\section{Summarium}

Es ist eine schöne Historia, wie Gott wunderbarer Weise schaffet, daß die Moabiter geschlagen werden. Solcher Sieg widerfäret dem Könige Israel und Edom allein um des frommen Josaphats willen.

Daß sie aber mit den Feinden so unbarmherzig umgehen, die Bäume umhauen, die guten Acker mit Steinen voll werfen, V. 25., geschieht nicht aus tyrannischem Fürnehmen, sondern darum, daß der Prophet Elisa solches ihnen zuvor gesagt hatte, daß es der Wille Gottes wäre.

Und hie stehet man wieder ein Exempel, wie es den Aufrührern gehet. Denn die Moabiter haben solche Strafe damit verdienet, daß sie von dem Könige Israel sind abgefallen.

\section{Votum}

Regiere uns der treue GOtt, daß wir ihm und unserer Obrigkeit schuldigen und willigen Gehorsam leisten; und helfe, daß wir uns in aller Noth seines heiligen Worts trösten und in wahrem Troste leben und sterben, Amen.

\chapter{Das 4. Kapitel}

\section{Vorrede}

Der königliche Prophet David sagt im 34. Psalm V. 20.: Der Gerechte muß viel leiden, aber der HErr hilft ihm aus dem allen. Das Kreuz ist die rechte Hof-Farbe den Christen, darin Christus seine Hofdiener kleidet. Gott salzet uns mit dem beißenden Salz des Kreuzes ein, daß nicht höllische Maden und verderbliche Sünden in uns wachsen und regieren. Gott höret uns im Ofen des Elends, daß wir nicht sinkend und faul, und demnach als untüchtig für die höllischen Hunde weggeworfen werden. Aber Gott hilft auch endlich von allem Kreuz, der da ist der einige Nothhelfer zur rechten Zeit, und der einige Meister zu helfen, Esaia 63, 1. Der alle Hülfe thut, die auf Erden geschieht, Ps. 47, 12. Weil dem also, daß uns alles Kreuz gut und nützlich ist, und daß uns Gott schuldlos will halten und endlich erretten und helfen, so sollen wir das Joch des HErrn Christi auf uns nehmen, unsere Seele mit Geduld fassen und auf Gottes Hülfe beständig warten. Davon hört auch an den ersten Theil des 4. Kapitels.

\section{Elisas fünf Wunderwercke}
\quellseite{548}

a (Ich wohne) das ist, ich habe zu Hofe nichts zu schaffen; ich wohne unter den Leuten alleine.

b Id est, More anni currentis, quo vivitur à cunctis animantibus. D. i. im Lauf dieses Jahres.

\section{Summarium}

Dies ist erstlich ein fein Exempel von einer armen Wittwe, in welchem man siehet, wie Gott, ob er gleich fromme Leute eine Zeitlang lässet Mangel leiden, dennoch endlich ihnen hilft und nicht in der Noth stecken läßt, V. 2.

Weil aber Gott der Wittwen hie solchen Segen läßt widerfahren, und heißet von dem Uebrigen die Schuldiger bezahlen, V. 7., siehet man, daß es Gottes Wille ist, daß, wer schuldig ist, dasselbe ehrbarlich und treulich bezahlen soll. Denn fremdes Gut ist durchaus verboten, daß es niemand soll inne haben. Daß aber Christus im Evangelio sagt, man soll leihen und nichts dafür nehmen, solcher Befehl gehet nicht dahin, daß man entlehnen und nichts soll wiedergeben, sondern daß die, so zu leihen haben, solches ohne ihren Nutzen und dem Dürftigen zu gut thun sollen. Denn das Entlehnte ist nicht eigen, und wo man kann, soll es wieder erlegt werden.

Zum andern ist ein schön Exempel von der reichen Frau zu Sunem, welche den Propheten Elisa beherberget und von Gott für solche Wohlthat einen Sohn erlanget. Denn Gott will sonderlich vergelten, was man denen zu gut thut, so sein Wort predigen.

\section{Votum}

Der treue, barmherzige GOtt helfe auch uns in aller Noth, und allen, die in Armuth stecken und andere Haus-Kreuze tragen; und regiere uns, daß wir auch Gerechtigkeit und Barmherzigkeit gegen unsern Nächsten üben, ihm zum Gehorsam und ihnen zu Hülfe und Trost, Amen.

\section{Der andere Theil des 4. Capitels}

\section{Vorrede}

Weil der Artikel von Auferstehung des Fleisches alle Furcht des Todes benimmt, und der heil. Geist dadurch lebendigen Trost in allen Anfechtungen in uns wirket; so bestätiget Gott diesen Artikel so reichlich mit Sprüchen und Exempeln des Alten und Neuen Testaments. Auch hat er derwegen denselben in vielen Creaturen abgebildet, als: wenn die Schnecken alt werden und fast verweset sein, werfen sie die alten Häute ab und werden wiederum aufs neue lebendig und jung. Die Seidenwürmlein werden auch wieder lebendig, ob sie schon eine Zeitlang todt gewesen sind. Alle Speise wird in das schönste rothe Blut verwandelt. Ein kranker Mensch wird gleich verzehret, bleich und erdenfarbig; aber wenn er wieder gesund wird, so wird er gleich wieder erquicket, fett und wohlgestaltet. Summa, wir sehen unterwärts, aufwärts, neben uns oder um uns, und sehen an Laub und Gras, Himmel und Erde, Menschen und Vieh, so sind sie unsere Haus-Postillen, die uns von der Auferstehung des Fleisches predigen. Gemeldeter Exempel eines erzählet auch der andere Theil des 4. Capitels.
\quellseite{549}

\section{Summarium}

Ist erstlich ein trefflich Exempel eines großen Glaubens an der Sunamitin, daß sie die Hoffnung hat, Gott werde ihren Sohn durch den Propheten vom Tod wieder zum Leben auferwecken, wie dann geschiehet.

Darnach sind zwei andere Exempel. Das erste, daß der Prophet Elisa die wilden Ranken, die vor Bitterkeit niemand essen konnte, wunderbarlich ändert, daß sie wohl schmeckten, V. 39. Das andere, daß er in der Theurung mit wenig Brod hundert Mann speiset, und noch übrig bleibt, V. 42. Solche Historien zeigen an, daß dennoch etliche fromme Leute im abgöttischen Königreich Israel sind gewesen, die an Gottes Wort gehalten haben, welchen er in der Theurung und sonst in andern Nöthen, um seines Wortes willen, geholfen hat; da die Ungläubigen haben müssen um ihrer Sünden willen Noth leiden. Uns zum Exempel, daß wir fromm sollen sein und hoffen, Gott werde es uns genießen lassen; gleich wie er die ungehorsamen bösen Buben mit der Strafe heimsuchen und ihrer auch nicht vergessen will.

\section{Votum}

Gott, der HErr, verleihe uns allen rechtschaffene Buße, erhalte uns in wahrem Glauben, speise uns an Leib und Seele, und sei unsere Hülfe in Noth und Tod, und erwecke uns auch zum ewigen Leben am jüngsten Tage, durch Jesum Christum, seinen lieben Sohn, unsern HErrn, Amen.
\quellseite{550}

\chapter{Das 5. Kapitel}

\section{Vorrede}

St. Paulus sagt Röm. 15, 9.: Daß die Heiden Gott loben um der Barmherzigkeit willen. Und Esaias spricht im 11. Kapitel V. 10.: Es wird sein die Wurzel Jesse, und der auferstehen wird zu herrschen über die Heiden, auf den werden die Heiden hoffen. Derer Sprüche ist ein Exempel Naeman, der Syrer. Ferner Salomo in Sprüchen im 21. Kapitel V. 6. sagt: Wer Schätze sammelt mit Lügen, der wird fehlen und fallen unter die, so seinen Tod suchen. Und Sirach spricht im 31. Kapitel V. 6.: Viele kommen zu Unfall um Geldes willen, und verderben drüber vor ihren Augen. Diese Sprüche sind erfüllet an Gehasi, dem Knecht Elisa. Die Exempel aber der erzählten Sprüche lauten aus dem 5. Kapitel.

\section{Naeman wird vom Aussatz gereiniget. Gehasi damit gestraft}

a Hand lehnet) das ist hebräisch geredet, wie wir Deutschen sagen: Er ist mir zur Hand, das ist, er ist um mich, thut und richtet aus, was ich ihm befehle und ich mich auf ihn verlasse.
\quellseite{551}

\section{Summarium}

Ist eine feine Historia, wie Naeman, der Syrer von seinem Aussatz rein und zum rechten Gott Israel bekehret wird, V. 14. Durch solche Historien hat Gott immerdar wollen anzeigen, daß auch die Heiden noch sollten zum Reich Gottes kommen und aus Gnaden ohne Gesetz selig werden.

Sonderlich aber ist hie zu merken, daß Gott seine großen Wunderwerke so gar einfältig und schlecht fürnimmt, daß sich die Vernunft muß darüber ärgern. Es gehe aber so schlecht und einfältig zu, als es immer wolle, wer Gottes Zusagung und Gnade gläubet und folget, der erfähret, daß sein Thun gewaltig und wunderbarlich hinaus gehet. Der Vernunft nach war der Jordan nichts besser, denn andere Wasser. Aber weil das Wort auf den Jordan weiset und auf andere Wasser nicht, wird Naeman rein, da er sich im Jordan badet. Es ist aber diese Historia auch ein Vorbild und Deutung der seligen Taufe, die im Jordan von Johannes angefangen und auf uns Heiden auch kommen ist, dadurch wir vom greulichen Aussatz der Sünden rein werden, durch den Glauben an den Tod und Verdienst Christi, wie es der heilige Basilius auch deutet.

b Wandelte nicht? das ist, hast du nirgend hin gewandelt; wie gehest denn zu, daß mein Herz wandelt, und war bei dem Wagen.

\section{Summarium}

Ist eine Historia, in welcher man siehet, wie Gott den Geiz und Lügen an des Propheten Knecht Gehasi straft, daß er und alle seine Nachkommen um solcher Sünden willen mit dem Aussatz gestraft werden, denn Gott ist allem Geiz feind. Sonderlich aber ist er dem Judas-Geiz feind, da man mit seinem heiligen Wort und Sacrament handthieret, daß man Geld, Gut, Herrschaft, Gewalt etc. dadurch erlange, wie der Erz-Judas, der Pabst, mit seinem Haufen thut, der aus Gottes Wort und den heiligen Sacramenten schlecht einen Jahrmarkt gemacht und sie um Geld verkauft hat.

\section{Votum}

Der allmächtige, barmherzige GOtt helfe heute auch an Leib und Seele allen Kranken und Kreuzträgern; und verleihe, daß wir uns für Geiz, Lügen und allen andern Sünden hüten, stürze und vertilge auch allen Judengeiz, Amen.

\chapter{Das 6. Kapitel}

\section{Vorrede}

Der heiligen Engel Schutz und das liebe tägliche Brod sind auch zwei große gnadenreiche Wohlthaten Gottes, unsers HErrn. Denn wo die heiligen Engel Hüter und Wächter sind, da können die Teufel und böse Menschen nicht schaden, weder mit List, noch mit Gewalt. Denn sie sind dienstbare Geister, die ausgesandt sind zum Dienste um derer willen, die ererben sollen die Seligkeit, Ebr. 1, 14. Und wo Gott das tägliche Brod entzeucht, da folget größere Noth, denn daß auszusprechen ist. Denn da müssen die barmherzigsten Weiber ihre Kinder selbst kochen, daß sie zu essen haben; wie zu Jerusalem geschehen, da diese Stadt verstöret war für der babylonischen Gefängniß, als davon Jeremias in Klagliedern im 4. Kapitel V. 10. sagt. Und in der letzten Zerstörung Jerusalems hat man Tauben-Mist für köstliche Würze gebraucht und theuer gekauft, und die Menschen haben ihren eigenen Mist gegessen, den Hunger zu stillen. Derwegen danken wir Gott heute herzlich, daß er uns durch seine heiligen Engel bisher beschü-
\quellseite{552}
het und für Theurung und Hunger behütet hat; der erbarme sich ferner unser aller gnädiglich, und gebe uns unsere Nothdurft und behüt uns für allem, was uns schädlich und verderblich sein mag, Amen. Dieses erinnert auch das 6. Kapitel.

\section[Vom schwimmenden Eisen. Blindheit der Syrer. Hungersnoth zu]{Vom schwimmenden Eisen. Blindheit der Syrer. Hungersnoth zu Samaria}

\section{Summarium}

Es ist eine sonderliche und schöne Historia, in welcher wir sehen, wie Gott um die, so ihn fürchten und seinem Wort anhangen, wachen, sie durch seine Engel beschützen und vor allem Uebel bewahren wolle, wider alles, was der Teufel und die Welt ihnen begehret zu thun. V. 17., wie der 34. Psalm V. 8. auch sagt: Der Engel des HErrn lagert sich um die her, so ihn fürchten, und hilft ihnen aus. Und der 91. Psalm V. 11.: Er hat seinen Engeln befohlen über dir, daß sie dich behüten auf allen deinen Wegen, daß sie dich auf den Händen tragen und du dei-
\quellseite{553}
nen Fuß nicht an einen Stein stößet etc. Durch solche Lehre und Exempel werden wir vermahnet, daß wir glauben sollen. (Denn mit den Augen werden wirs eben so wenig sehen, als des Propheten Knecht hie,) daß wir mit solchem Schutz und Hut der lieben Engel ohne Unterlaß umgeben und bewahret sind, sonst wären wir Leibes und Lebens, und alles, was wir haben, nicht einen Augenblick für dem Teufel sicher.

a Uebel ja, so gehets uns, wenn wir unserm Gott dienen, ihr heillosen Propheten. Wie viel besser hatten wirs, da wir Baal dieneten, Jer. 44, 17.

\section{Summarium}

Ist eine Historia, wie Gott den abgöttischen König mit Krieg und großer Theurung strafet, daß die Mütter ihre eigenen Kinder erwürgeten, kochten und aßen. Solches Unglück pfleget Gott zu erwecken, wenn man sein Wort verachtet und sich von falscher Lehre und unrechtem Gottesdienst nicht abwendet. Denn eben um solcher Ursache willen hatte Gott diese großen Propheten, Elia und Elisa, erwecket, daß die Abgötterei abgethan und rechter Gottesdienst wieder würde angerichtet. Aber weil sie sich an ihr Predigen und Wunderzeichen gar nichts kehren, strafet sie Gott mit solcher großen Theurung.

\section{Votum}

Gott, unser aller lieber Vater, behüte uns gnädiglich für Krieg, Theurung und andern Strafen; und beschütze durch seine heiligen Engel dieses Land und Stadt, uns und die Unsern, und verleihe, daß wir nach seinem Wort all unser Thun und Lassen allzeit anstellen, Amen.

\chapter{Das 7. Kapitel}

\section{Vorrede}

Gleich wie Gott uns tröstet und erfreuet in unserer Noth; also sollen wir auch unserm Nächsten in seinem Nothstande tröstlich sein und ihn erfreuen. Denn das ist einmal wahr: Je größer Noth, je näher Gott. Und abermal ist recht gesagt: Unser Gott ist über alle Noth. Derwegen der nothleidende Mensch gar wohl bauet, wenn er fest auf Gott trauet. So lehret auch der Herr Christus dieses Matth. 7, 12.: Alles, das ihr wollet, daß euch die Leute thun sollen, das thut ihr ihnen. Und Tobias sagt im 4. Kapitel V. 16. zu seinem lieben Sohn auch also: Was du willst, daß man dir thue, das thue einem andern auch. Welche Worte unsere Richtschnur und Regel sein sollen, darnach wir alle unsere Worte, Werke, Gedanken und Gebärden gegen unsern Nächsten anstellen sollen, zu gefallen Gott und zu Nutz andern Leuten. Dies lehret auch das 7. Kapitel.

\section{Geschwinde wohlfeile Zeit zu Samaria}
\quellseite{554}

\section{Summarium}

Dies ist eine herrliche schöne Historia, wie Gott die Syrer wunderbarlich schrecket und jaget, v. 6., und Samaria von dem jämmerlichen Hunger erlediget und nach der größten Theurung plötzlich ein unerhörter und unglaublicher Ueberfluß aller Dinge erfolget, v. 10.

Diese und dergleichen Historien dienen, daß wir an Gottes Wort fest halten und in der Noth seiner Hülfe uns gewiß trösten sollen, als deß,
\quellseite{555}
der allmächtig ist und tausenderlei Wege zu helfen hat, da wir in solcher Noth nicht einen wissen noch sehen. Auf daß wir nicht thun, wie der gottlose König und seine Ritter, die dem Propheten nicht glauben wollen. Auch ist hie zu merken das Exempel der Aussätzigen, daß sie es dafür halten, es wäre Sünde, wenn sie allein ihren eigenen Nutzen bedenken und den armen bedrängten Leuten nicht bald solche gute Botschaften verkündigen sollten.

\section{Votum}

Gott, der heilige Geist, regiere uns, daß all unser Thun und Lassen unserm Nächsten auch zur Freude und Nutz diene; und tröste und helfe allen, die in Noth und Gefahr stecken, damit sie seine Gnade und Barmherzigkeit zeitlich und ewig rühmen und preisen, Amen.

\chapter{Das 8. Kapitel}

\section{Vorrede}

Gott ist gerecht und auch gnädig, zornig und auch barmherzig. Gerecht ist er gegen die Gottlosen, die in Sünden verstarren und verharren. Gnädig erzeigt er sich gegen die Frommen, die Tugend lieben und üben. Zornig wird Gott um der Sünde willen; barmherzig ist er, wenn man die Sünde erkennt und die ihm abbittet im Namen Jesu Christi, seines lieben Sohns, des unschuldigen Lammes Gottes, das der Welt Sünde getragen und gebüßet hat, und vergibt die Sünde, wendet auch ab oder lindert die wohlverdienten Sünden-Strafen. Derwegen laßt uns in wahrer christlicher beständiger Buße leben, daß uns Gott in seinem Zorn nicht strafe, sondern aus Gnade an Leib und Seele segne und behüte, gebe und erhalte auch, was uns noth ist. Das helfe uns Gott allen, Amen. Darauf höret auch an das 8. Kapitel.

\section[Siebenjährige Theurung. Veränderung im Königreich der Syrer. Joram]{Siebenjährige Theurung. Veränderung im Königreich der Syrer. Joram und Ahasja, Könige in Juda}
\quellseite{556}

\section{Summarium}

Ist ein Exempel, daß die Sunamitin wiederum deß genießet, daß sie so gastfrei gewesen, und dem Propheten Elisa so viel Gutes gethan hatte. Denn der Prophet warnet sie, daß sie sich aus dem Lande machen und der Theurung entfliehen sollte, V. 1. Also will Gott mit reichem Segen bezahlen, was man seinem Worte zu Ehren denen Gutes thut, die es predigen und lieb haben. Darnach ist es ein Exempel, wie Gott den Tyrannen Benhadad sterben läßt, V. 10. Aber weil die Kinder Israel sich nicht bessern, kömmt Hasael, noch ein ärgerer König, an seine Statt, der ihnen mehr zu Leide that, denn Benhadad.

\section{Summarium}

Ist ein Exempel, wie der fromme König Josaphat einen ungerathenen Sohn hatte, mit Namen Joram. Darum plaget denselben Joram Gott erstlich mit Krieg, daß die Edomiter von ihm abfallen, V. 20. Darnach plagt ihn Gott mit schwerer Krankheit, der man nicht rathen konnte, daran er auch starb, da er zwei ganze Jahre mit Qual darnach lag. Davon weiter 2 Chron. 21, 8. und 14. 15. 16. Dergleichen Ahazia, sein Sohn, weil er gottlos ist, regieret auch nur ein Jahr und stirbt, V. 29. Also siehet man, daß Gott weder Glück noch Heil den Regenten gibt, so gottlos sind und Gottes Wort verachten.

\section{Votum}

Der gerechte GOtt steure auch heute allen muthwilligen Verächtern seines Wortes; und vergelte zeitlich und ewig alles, was man Gutes thut seinen Dienern, und behüte uns für Krieg, Krankheit, Theurung und anderm Elend und Herzeleid, Amen.

\chapter{Das 9. Kapitel}

\section{Vorrede}

Der heilige Prophet Jesaias sagt im 8. Kapitel V. 20.: Nach dem Gesetz und Zeugniß. Werden sie das nicht sagen, so werden sie die Morgenröthe nicht haben. Das ist wahr geworden an König Joram und Ahazia. Ferner sagt der Prophet Zacharia im 2. Kapitel V. 8.: Wer euch (vernehmt die treuen Prediger)
\quellseite{557}
antastet, der tastet Gottes Augapfel an. Dies ist erfüllet an der Königin Isebel. Auch ist's gewiß, die den HErrn Christum (die rechte Morgenröthe, den Aufgang aus der Höhe, die Sonne der Gerechtigkeit, und den ewigen Glanz der ewigen Herrlichkeit Gottes, des Vaters) und sein Wort verachten und verwerfen, die sollen hie zeitlich gestraft werden mit Hunger, Krieg, Theurung, Krankheiten und andern Plagen, sie sollen auch hernach in den Abgrund der Hölle verworfen werden, zu ewiger Marter und Pein. Dieses werden bezeugen alle Gottlose in ewiger Verdammniß. Die Geschichte gemeldeter Straßen und Untergangs lautet aus dem 9. Kapitel.

\section[Jehu zum Könige gesalbet, tödtet Joram, Ahazia und]{Jehu zum Könige gesalbet, tödtet Joram, Ahazia und Isebel}

a Rasender) Non quod attoniti aut devoti, ut Rabini delirant, sed quod impii Prophetas vocant furiosos, sicut hodie. D. i. nicht daß sie von Sinnen oder schwärmend gewesen, wie die Rabbinen träumen; sondern weil die Gottlosen die Propheten Wahnsinnige heißen, wie noch heute.

b Hohe Stufen) hic fingendum est, fuisse Sellam Magistratus elevatam in urbe, candidam quasi osseam: huc posuerunt Jehu. Sed quia pompa regia tam cito non poterat tapetis ornari, suas vestes substernebant, in pompa festinantes scilicet. D. i. hier muß man sich vorstellen, es sei ein erhöhter Amts- oder Gerichtstuhl in der Stadt gewesen, weiß wie elfenbein: auf denselben setzten sie den Jehu. Aber weil der königliche Pomp nicht so schnell konnte mit Tapeten oder Teppichen hergerichtet werden, so breiteten sie ihre Kleider unter; weil sie nämlich bei diesem Pompe eilten.
\quellseite{558}

\section{Summarium}

Ist eine Historia, wie der gottlose König in Israel, mit Namen Joram, sammt dem König Ahazia in Juda, aus Befehl Gottes um ihrer Abgötterei willen erwürget werden, V. 24. 27. Auf daß man ja lerne und es gewiß dafür halte, wer sich an Gottes Wort nicht kehren, dasselbige verachten und von der Abgötterei und falscher Lehre nicht ablassen will, daß derselbige gewiß gestraft müsse werden. Derhalben mag sonderlich die weltliche Obrigkeit, wo sie anders Glück und Heil haben will, sich an Gottes Wort begeben und demselbigen mit Herzen anhangen und allem abgöttischen Wesen mit Ernst wehren.

\section{Summarium}

Hie ist eine sonderliche Historia, wie es mit der gottlosen Isebel habe ein Ende genommen; zum Exempel allen denen, so Gottes Wort feind sind, und die, so es predigen, verfolgen; daß es Gott ihnen nicht schenken, sondern sie gewißlich strafen will, und soll sie weder ihre Gewalt noch Reichthum helfen. Isebel war eine Königin und eines Königs Tochter, aber mit ihrem gottlosen Tyrannenfürsten, wider Gott und sein Wort.

33 *
\quellseite{559}
verdienet sie, daß sie zum Fenster heraus geworfen und von den Hunden gefressen wird. Also verachtet und schändet Gott wiederum die, so ihn und sein Wort verachten und dawider toben. Geschichts nicht auf solche Weise, wie hie mit Isebel, so geschicht es doch auf andere, wie wir zu unsern Zeiten viele Exempel erlebt haben.

\section{Votum}

Der gerechte GOtt wolle noch steuren und wehren den Feinden seiner Kirche, und erhalten, die sein Wort recht und treulich lehren, und die seinen heiligen Namen bekennen, loben und preisen; und wolle uns regieren, daß wir uns für allen Sünden, wider ihn, unsere Lehrer und Obrigkeit, allzeit fleißig hüten, Amen.

\chapter{Das 10. Kapitel}

\section{Vorrede}

Höret an mit andächtigen Herzen die Worte unsers Gottes und Herrn, da er im 2. Buch Mose im 20. Kapitel V. 5. also sagt: Ich, der Herr, dein Gott, bin ein eifriger Gott, der da heimsuchet der Väter Missethat an den Kindern, bis ins dritte und vierte Glied, die mich hassen; und thue Barmherzigkeit an viel Tausenden, die mich lieb haben und meine Gebote halten. In diesen Worten dräuet Gott zu strafen alle, die seine Gebote übertreten, beides an ihrer Person und ihren Nachkommen. Darum sollen wir uns fürchten vor seinem Zorn und nicht wider seine Gebote thun; er verheißet aber auch Gnade und alles Gute allen, die seine Gebote halten, beides an ihrer eigenen Person und auch an ihren Nachkommen. Darum sollen wir ihn auch lieben und vertrauen, und gerne thun nach seinen Geboten. Davon unterrichtet auch das 10. Kapitel des 2. Buchs der Könige in den Exempeln Ahab und Jehu.

\section{Ausrottung des Geschlechts Ahabs und der Baals-Pfaffen}
\quellseite{560}

\section{Summarium}

Ist eine Historia, in welcher man die Strafe siehet, welche Gott über den abgöttischen König Ahab hat kommen lassen, daß er nicht allein, sondern auch seine Kinder und Verwandten vertilget müssen werden, V. 7. Solche Exempel sollten uns eine Warnung sein, daß wir uns für falscher Lehre und unrechtem Gottesdienst mit Ernst hüten und der gottseligen Lehre anhangen, unangesehen, daß man etwas drüber wagen muß; denn es ist besser, einen gnädigen Gott und die Welt zum Feinde haben, denn der Leute Gunst und Gottes Zorn und Ungnade tragen.

a Das ist über die Cammer.
\quellseite{561}

\section{Summarium}

Jst ein Exempel einer sonderlichen und fürstlichen Tugend, die Gott wohl gefällt und mit reichem Segen belohnet, daß Jehu alle falsche Propheten, so zum abgöttischen Dienst des Baals geholfen hatten, erwürgen lässet und die Kirche Baal zum heimlichen Gemach machet. Denn weltlicher Obrigkeit gehöret dies sonderlich zu, daß sie mit allem Ernst und ohne alle Gnade der Abgötterei wehre. Denn weil es Gottes Ehre und des Nächsten Seligkeit belanget, soll man darin nicht faul noch lässig sein. Das will Gott hernach mit glückseliger Regierung bezahlen. Wie Jehu ganzer acht und zwanzig Jahr im Regiment bleibt und sein Stamm bis ins vierte Glied. Eines solchen Jehu bedürfte heutiges Tages sehr wohl das Papstthum, sammt allen Stiftern und Klöstern, darin eben so große und greuliche Abgötterei getrieben wird, als dort von den Baalspfaffen.

\section{Votum}

Der eifrige gütige Gott wehre heute auch aller Abgötterei und allen Götzendienern, und segne an Leib und Seele, sammt allen ihren Nachkommen, alle, die ihn lieben und fürchten, und thue solches um Jesu Christi, seines lieben Sohnes, willen, Amen.

\chapter{Das 11. Kapitel}

\section{Vorrede}

Gottlose Obrigkeit ist ein allgemeiner Landschaden; deß ist ein Exempel Athalia. Dagegen gottesfürchtige Obrigkeit ist ein allgemeiner Landfrommen; das lehret das Exempel Jojada. Demwegen danken wir heute aber Gott, daß er durch unsere liebe Obrigkeit sein heilig Wort erhält, und daß durch sie die Unterthanen väterlich regiert und treulich geschützt werden. Der wolle ferner zum Regiment dieser Stadt und Landes geben Glück und Segen, Gesundheit und Leben, beides dem hochlöblichen römischen Adler und dem Rathstuhl der Stadt Obrigkeit, Amen. Ferner höret an das 11. Kapitel.

\section{Tyrannei Athaliae. Das Krönung}

a Des Sabbaths) das waren, die auf des Königes Dienst warten, eine Woche um die andere. Wenn ein Theil abging, so ging das andere an.

b Zeugniß) das war das Buch Mose, das fünfte, das dem König befohlen war.
\quellseite{562}

\section{Summarium}

Dies ist eine merkliche Historia, wie der Priester Jojada Joas, des Königs Ahazia Sohn, zum Könige machet, V. 12., und Athalia, welche das Regiment an sich gezogen und die Kinder des Königs alle erwürget hatte, umbringen heißt, zum Exempel allen Tyrannen, daß sie sich für Tyrannei hüten oder dergleichen Strafe auch warten sollen.

Darnach ist wieder ein Exempel, wie der Priester Jojada die Kirche Baals einreißet, Baals Priester erwürget und den rechten Gottesdienst, so unter dem Regiment Athalia gefallen war, wieder anrichtet, V. 18. Solche Werke gehören sonderlich den Regenten zu, so anders Glück und Heil bei ihrem Regiment sein soll.

\section{Votum}

Der allmächtige Gott steure auch heute allen Tyrannen, und regiere und erhalte unsere treue Obrigkeit zu seiner Ehre und zu gemeinem Nutz, und helfe, daß wir derselben allezeit schuldige Ehre und christlichen Gehorsam erzeigen, Amen.

\chapter{Das 12. Kapitel}

\section{Vorrede}

Vom Amt christlicher Obrigkeit sagt der königliche Prophet David im 24. Psalm V. 7.: Machet die Thore weit und die Thüren in der Welt hoch, daß der König der Ehren einziehe. Mit dem Wort Thor verstehet David die Regimente und Regenten, weil man bei den Juden in Thoren Gericht pfleget zu halten; und vermahnet alle hohen und niedrigen Standes Regenten, daß sie nicht allein darauf bedacht sein sollen, daß sie in ihrer Jurisdiction und Gebiet gute äußerliche Policei, Friede und Ruhe erhalten, sondern daß auch Christus, der Fürst des Lebens, Apost. 3, 15., und der Seligkeit, bekannt und aufgenommen werde; daß sie also für ihre Person die Wahrheit und Klarheit, die Ehre und Lehre Christi, lieb und werth halten, dieselbige bewahren und beherbergen, und daß sie auch der christlichen Kirche, Patronen, Schutzherren und Pfleger sein sollen und die ganze Religion fördern und erhalten. Denn das ist der rechte Ornat, Zierrath und Schmuck aller Regimente, wie Kaiser Theodosius gesagt hat. Ja das ist die Thür, dadurch der liebe Gott mit seinem Segen himmlischer und irdischer Wohlfahrt einzeucht bei Obrigkeit und Unterthanen. So helfe der wahre gütige Gott, daß alle Obrigkeit ihr Amt also löblich und glücklich verrichte, wie auch gethan haben Jojada und Joas, von denen das 12. Capitel.

\section{Regierung Joas, Königs in Juda}
\quellseite{563}

\section{Summarium}

Sie ist eine Historia, wie Jojada sammt dem König ordnen, daß der Tempel im Bau erhalten werde, V. 5. Da siehet man wieder, wie es weltlicher Obrigkeit Amts halben gebühren will, daß sie sich mit Ernst der Religion, und was dazu gehöret, sollen annehmen. Sonderlich aber ist zu merken, daß der Text meldet: Joas that, was recht war und dem HErrn wohl gefiel, so lange ihn der Priester Jojada lehret, V. 2. Denn wenn weltliche Obrigkeit, soll sie anders allenthalben recht fahren, muß sie die an der Hand haben, ihres Raths und Unterrichts oft brauchen, welche Gottes Wort predigen und damit umgehen. Darum stehets nimmermehr wohl, wo man die Prediger und Kirchendiener verachtet, wie doch gemeiniglich jetzund die Obrigkeit allenthalben thut.

\section{Votum}

Regiere GOtt ferner unsere und andere christliche Obrigkeiten, daß sie ihnen die Religion lassen befohlen sein, ihnen selbst zu gutem Gewissen und den Unterthanen zu Heil und Wohlfahrt, und segne und schütze sie an Leib, Seele, Land und Leuten, Amen.
\quellseite{564}

\chapter{Das 13. Kapitel}

\section{Vorrede}

Um der Abgötterei willen sind die Cananiter zu Grunde ausgerottet und die Juden verworfen und in alle Welt zerstreuet worden; also sind die mächtigsten und gewaltigsten Länder auf Erden, Syrien, Egypten, Griechenland und das ganze orientische Kaiserthum, auch um der Abgötterei willen, in der Türken Gewalt kommen. Wehe nun der unbußfertigen Welt der Abgötterei wegen, die heute im Schwange gehet und getrieben wird. Denn Gott ist der Herr, unser Gott, darum soll sich niemand zu den Götzen wenden, denn der Herr hasset alle Abgötterei und alle Abgötter müssen umkommen, 3 Mos. 19., Sir. 15., 5 Mos. 8. Derwegen wer beten kann, der bete alle Tage, christlich, herzlich und fleißig, daß Gott unser sich erbarme, und daß er uns fremder Abgötterei nicht entgelten lasse. Darauf höret auch an das 13. Kapitel.

\section{Regierung Joahas und Joas}
\quellseite{565}

\section{Summarium}

Dies ist eine Historia, wie Gott das Königreich Israel mit heftigen und großen Kriegen plaget, V. 3. Alles um der Abgötterei willen. Dennoch da sie um Hülfe bitten, schaffet ihnen Gott Hülfe durch den König Joahas, um seines Bundes willen mit Abraham, Isaak und Jakob, das ist, aus Gnade und Barmherzigkeit.

Die Historia mit dem Grab Elisa, des Propheten, V. 21., muß man deuten, wie andere seine Wunderwerke, auf daß das Volk sähe und wüßte, daß der Gott, von dem Elisa geprediget, der rechte Gott und seine Lehre die rechte Lehre wäre, und sie desto eher von der Abgötterei ließen. Dagegen die Papisten die erlogenen Wunderwerke ihrer Heiligen allein zur Bekräftigung der Abgötterei wider das Wort gebraucht haben, und wollens noch heutiges Tages mit diesen und andern Wunderwerken schmücken. Aber es reimet sich sehr übel, wie gehöret.

\section{Votum}

Der treue GOtt erleuchte alle, die in Abgötterei stecken, daß sie die abschaffen und ihm im Geist und in der Wahrheit dienen; und behüte uns für Krieg und allem Unglück, tröste auch heute alle fromme Christen in aller Noth, Amen.

\chapter{Das 14. Kapitel}

\section{Vorrede}

Laßt uns heute abermal Gott, dem Allmächtigen für die Schenkung seines Worts und Offenbarung seines Willens danken; und ein jedes kehre auch allen möglichen Fleiß an, daß wir dieser Gnade Gottes nützlich gebrauchen, nach der Vermahnung des heiligen Propheten Esaia am 55. Kapitel V. 6., da er spricht: Suchet den HErrn, weil er zu finden ist, rufet ihn an, weil er nahe ist, der Gottlose lasse von seinem Wege, und der Uebelthäter seinen Gedanken, und bekehre sich zum HErrn, so wird er sich seiner erbarmen, und zu unserm Gott, denn bei ihm ist viel Vergebung. Den rechten Weg zu glauben, und solches zur Seligkeit, den zeiget uns an das Evangelium von Jesu Christo; den Weg recht zu leben, und dies zu unserer Wohlfahrt, den zeigen uns die heiligen zehn Gebote Gottes, dazu unsere christliche Haus-Tafel und eines jeden Beruf und Stand. Diese Lehre hat Gott bestätiget mit Exempeln, beide der Belohnung und Strafe an denen, die darnach oder dawider gelebet haben. Hierher gehört auch das 14. Kapitel.

\section[Amazia und Asaria, Könige in Juda; Jerobeam der]{Amazia und Asaria, Könige in Juda; Jerobeam der andere in Israel}
\quellseite{566}

\section{Summarium}

Ist erstlich eine Historia, wie der Mord, welchen die Unterthanen an ihrem Könige Joas begangen hatten, gerochen wird. Und wird sonderlich an Amazia gelobet, daß er solche Strafe an den Thäter wendet und ihre Kinder, sammt der andern Freundschaft, die daran nicht schuldig waren, solches nicht entgelten lässet. Darnach ist eine andere Historia, wie Amazia einen Krieg wider das Königreich fürnimmt und darüber geschlagen und ein langer Ort der Mauern zu Jerusalem eingerissen und die Stadt geplündert wird, V. 13. Alles darum, daß die Sünde Amazia und Juda gestraft würde, daß sie die Götter der Edomiter gesucht und der König des Propheten Strafe und Warnung verachtet hatte, und half sie nichts, daß Israel auch abgöttisch und voll Sünde und Laster war.

\section{Summarium}

Ist eine feine Historia, wie Amazia endlich erschlagen wird, darum, daß er sich an die Abgötterei gehänget und Gottes Wort, so durch den Propheten ihm gepredigt, verachtet hatte.
\quellseite{567}
Die Historia von Jerobeam, dem Könige Israel, von dem die Schrift zeuget, daß er auch habe gethan, daß dem HErrn übel gefiel, und doch 41 Jahr regierte und viel Glück mit Kriegen wider die Syrer gehabt, V. 26. Ist darum wohl zu merken, daß man lerne, daß Gott so gnädig und gütig sei, daß er sich des Elendes, auch derer, so die Strafe wohl verdienet haben, jammern läßt, ohne Zweifel um weniger frommer Leute willen, so noch vorhanden waren, und sich an der rechten Propheten Lehre gehalten, und ihnen gefolget hatten, derselben genießet der König Jerobeam sammt dem ganzen Haufen.

\section{Votum}

GOtt, der da gnädig und barmherzig ist, erhalte uns auch allzeit einen heiligen Samen und Leute, die ihm dienen in Heiligkeit und Gerechtigkeit, und lindere uns die Strafen, die wir mit so viel und großen Sünden verdienet haben; und helfe, daß die Obrigkeit das Böse strafe und das Gute belohne, Amen.

\chapter{Das 15. Kapitel}

\section{Vorrede}

Gleich wie drei große Krebse und Staubbesen sind, als Krieg, Theurung, Pestilenz; also sind auch drei große Hauptsünden, Abgötterei, Tyrannei und Unzucht; aber Abgötterei ist die größte Sünde und wird auch am heftigsten gestraft. Man soll aber wohl lernen, was Abgötterei sei, nämlich: Abgötterei ist zum ersten, dasjenige zu Gott machen, was nicht Gott ist. Zum andern, die Ehre, so Gott allein gebühret, einem Dinge geben, das nicht Gott ist. Zum dritten, Gott an einen Ort oder Kreatur verbinden, dahin er sich nicht verbunden hat in seinem Wort. Zum vierten, vermeinte Gottesdienste dichten, erwählen und verordnen ohne Gottes Wort und Befehl, und glauben, daß man um derselben willen Gott gefalle, gerecht und selig werde. Diese greuliche und verdammte Sünde der Abgötterei laßt uns ja fleißig fliehen, auf daß wir mit den Abgöttern nicht zeitlich und ewig gestraft werden; wie solcher Strafe Exempel uns aber vorgehalten werden im 15. Kapitel.

\section{Regierung etlicher Könige in Juda und Israel}
\quellseite{568}

\section{Summarium}

Ist eine Historia, in welcher man siehet, wie Gott an beiden Königen, Israel und Juda, die Abgötterei strafet, nicht allein mit Krieg und fremdem Volk, sondern auch mit Aufruhr, daß im Königreich Israel immer einer den andern erwürget, V. 10. 14. 19. 25.

Sonderlich aber ist zu merken das Exempel mit dem Könige Asaria, welcher sonst rechtschaffen und fromm war; aber darum, daß er die Höhen nicht abthat, und das Volk nicht dazu hielt, daß sie allein zu Jerusalem opferten, da der HErr befohlen hatte zu opfern, strafet ihn Gott mit dem Aussatz, V. 5. So wird auch sein Sohn Jotham überfallen von den Syrern, V. 37. Auf daß man lerne, daß sich die welt-


\quellseite{569}
liche Obrigkeit schwerlich damit versündiget, wenn sie der Abgötterei nicht mit Ernst wehret, ob sie gleich selbst nicht Abgötterei treibet.

\section{Votum}

Wir danken GOtt, daß er uns sein Wort reichlich und richtig gegeben und von mancherlei Abgötterei gnädiglich entlediget hat; und bitten ihn herzlich, er wolle uns mit seinem Geist regieren und führen, daß wir die Zeit unserer Heimsuchung erkennen, dankbar sein und für allen Sünden uns hüten, Amen.

\chapter{Das 16. Kapitel}

\section{Vorrede}

Die Sünden und Sündenstrafen der Gottlosen sind uns in der heiligen Bibel zur Warnung vorgeschrieben: die Sünden, daß wir sie hassen und lassen sollen; die Strafen, damit wir uns nicht zeitliches und ewiges Verderben auf den Hals ziehen, zu Leibes- und Seelen-Schaden. Derwegen helfe Gott, der Brunnquell und Geber alles Guten, daß wir sein Wort also anhören, daß die Sünder sich bekehren, daß die Frommen in Gottesfurcht beständig verharren, daß die betrübten Herzen getröstet und daß wir alle Gottes Segen zeitlich und ewig haben und behalten. Und gedenket, als wenn Moses vor einem jeden jetzt stünde und ihn anspräche: So du wirst halten alle diese Gebote, die dir Gott geboten hat, daß du darnach thust, so wird aller Segen über dich kommen; wenn du aber nicht gehorchen wirst der Stimme des HErrn, deines Gottes, daß du haltest und thust alle seine Gebote und Rechte, so werden alle Flüche über dich kommen und dich treffen, 5 Mos. 28, 1. 15. Denn also können auch die Worte Mosis gebraucht werden zur Warnung. Folgendes höret an das 16. Capitel.

\section{Ahas und seine Regierung}
\quellseite{570}

a Decke) darunter sie des Sabbaths saßen oder stunden, wie jetzt Fürsten und Herren unter Teppichen oder Getäfelten sitzen.

\section{Summarium}

Ist eine Historia, wie der König Ahas Abgötterei getrieben und neue Gottesdienste aus eigenem Fürnehmen, ohne Gottes Befehl, habe angerichtet und darum vom Könige zu Syrien sei überzogen worden, V. 5. In solcher Noth suchet er nicht Hülfe bei Gott, sondern bei dem Könige zu Assyrien, V. 7. Durch denselben werden die Syrer und endlich auch das Königreich Israel gestraft, wie sie verdient hatten, V. 9. Auf daß man lerne fromm sein, Gott fürchten und nach seinem Wort sich halten, sonst würde dergleichen Unrath nicht außen bleiben.

Sonderlich aber ist das Exempel mit dem Priester Uria zu merken: der beuget dem Könige und thut in der Kirche, was ihn der König heißet; so doch das Widerspiel sollte sein, daß der König seines Regiments warten und der Priester in der Kirche nicht auf den König, sondern auf Gottes Wort sollte sehen, so ginge es zu beiden Theilen recht. Aber der Teufel siehet solche Ordnung nicht gern. Die Ordnung ist ihm viel dienstlicher, darum hilft er dazu, wie man mit des Papstes und Bischöfe Regimenten auch siehet, daß es durchaus ein umgekehrt Ding ist. Sie sollten Prediger und Lehrer sein, so wollen sie Regenten sein; Könige und Fürsten sollen Regenten sein, so wollen sie die Leute zwingen, sie sollen glauben, zum Sacrament gehen und anderes mehr thun, was ihnen gefällt. Solche Unordnung aber bringet gewissen Unrath mit sich.

\section{Votum}

Der treue GOtt lasse ihm das weltliche und geistliche Regiment in diesen letzten und gefährlichen Zeiten befohlen sein; und helfe, daß Obrigkeit und Prediger ihrem Amt recht und nützlich vorstehen, und daß Unterthanen und Zuhörer christlich und gehorsamlich folgen, Amen.

\chapter{Das 17. Kapitel}

\section{Vorrede}

Der heilige Prophet Esaiaß sagt am 30. Kapitel V. 12.: So spricht der Heilige in Israel: Weil ihr das Wort verwerfft und verlasset euch auf Frevel und Muthwillen und trotzet darauf; so soll euch solche Untugend sein, wie ein Riß an einer hohen Mauren, wenn es beginnet zu rieseln, die plötzlich unversehens einfällt und zerschmettert. Weh nun uns Deutschen; denn in Kirchen unter den Gelehrten ist groß Geschrei, Gift und Neid. Im weltlichen Regiment gehets auch seltsam, fast überall. Die Kinderzucht fällt hinweg. Jedermann beinahe wird sicher, roh und gottlos. Daher uns so viele Himmelszeichen und alle Elemente dräuen, und zeigen an, Gott habe freilich was Großes im Sinn und wolle die Ruthe und die verdienten Strafen gehen lassen. Derwegen, o HErr Gott, o Vater aller Gnaden und Barmherzigkeit, heile du uns, so werden wir heil, hilf du uns, so ist uns geholfen, Jer. 17, 14. Amen. Dies soll uns auch erinnern das 17. Kapitel.

\section[Unter Hosea werden die zehn Stämme nach Assyrien]{Unter Hosea werden die zehn Stämme nach Assyrien geführet}
\quellseite{571}

a Schmückten) sie wollten ihre Sünde vertheidigen als recht und wohlgethan, wie alle Ketzer und Abgöttische thun.

\section{Summarium}

Hie ist der fürnehmsten Historien eine, in welcher man sehet, was für ein Unglück und Jammer daraus komme, wenn man Gottes Wort verachtet, Abgötterei treibet und sich an seine Strafe nicht kehren, noch dadurch sich bessern will: nämlich, daß Gott das ganze Israel, die zehn Stämme, durch ein fremdes Volk gefangen wegführen und ihr Land verheeren läßt, V. 18. Wie wir zu unsern Zeiten leider auch sehen, und täglich, wie zu besorgen, je länger je mehr erfahren werden, daß der Türke mit uns dergleichen wird umgehen, weil des Lästerns und Verfolgens wider die heilige erkannte Wahrheit auf der Papisten Seite kein Ende ist, auf unserer Seite aber die Verachtung und der Undank gegen das liebe Wort so groß wird.

Darum soll man hie sonderlich sehen, wie der Text die Ursachen nach einander herzählet, so solche Strafe erregt haben, nämlich, daß sie fremde Götter angebetet, V. 7., die Warnung der Propheten nicht angenommen, sondern verachtet und ein jeglicher nach seinem Gutdünken ihm eine eigene Weise, Gott zu dienen, habe fürgenommen, V. 15.; wie leider zu unsern Zeiten auch geschieht.
\quellseite{572}



\section{Summarium}

Ist eine Historia, in welcher wir sehen, wie und aus was Ursachen die heidnische Abgötterei in Samaria erstlich gepflanzet und angerichtet sei, durch das fremde Volk, welches der König von Assyrien anstatt der Juden drein hat gesetzt. Solche Abgötterei strafet Gott auch an den Heiden, und schicket Löwen unter sie, welche ihrer viel umbracht(en), V. 24. Ob sie wohl nun das erkennen, daß solches eine Strafe Gottes sei, und sich darum unterrichten lassen vom rechten Gottesdienst, dennoch bleiben sie nicht dabei, sondern fallen wieder auf ihr altes abgöttisches Wesen, V. 28. 29. Denn der Teufel nimmt sich darum mit höchstem Fleiß und Ernst an, daß er die Leute vom Wort und rechten Gottesdienst abwenden möge. Darum bedarf es fleißiges Aufsehens und ernstliches Gebet, daß uns Gott dabei gnädig bis ans Ende erhalten wolle. Es ist aber eine greulichere Strafe auf solche Sünde gefolgt, denn auf die Leute zuvor gewesen war. Denn es ist auf die Abgötterei Blindheit und Verdammniß der Seelen gefolgt.

\section{Votum}

Gott erhalte uns bei seinem heiligen Wort und rechten Gottesdiensten, und lindere die wohlverdienten Strafen; helfe auch, daß wir nach seinem Gesetz wandeln, sein Zeugniß halten, ihn von ganzem Herzen suchen und kein Uebels thun, Amen.

\chapter{Das 18. Kapitel}

\section{Vorrede}

Unser lieber HErr Gott hat eine hebräische Sprache. Er lehret das hinterst zuförderst, erst das Leiden, darnach die Herrlichkeit, hat der heil. Lutherus pflegen zu sagen. Das gibt auch
\quellseite{573}
die Ordnung des Jahrs. Die Marterwoche gehet allzeit vor dem Oster-Tage her, und auf die Kreuzwoche das Pfingstfest erst folget. Wer auch in unsers Herrn Gottes Lustgärten (den er hier auf Erden hat, im Wort und Sacramenten) spazieren will, der denke nur nicht, daß er eitel Tausendschönchen darinnen finden werde, sondern er wird auch bisweilen Sauerampfer und bittere Wermuth antreffen, welche doch gleichwohl auch nicht ungesund, sondern gut ist, Leib Lust und Gift zu vertreiben, die muß man auch kosten und schmecken, sondern werden wir vom Teufel und eigenem Fleisch gar bald vergiftet und angesteckt werden. Davon lehret auch die Epistel an die Ebräer im 12. Kapitel V. 6. mit solchen Worten: Welchen der Herr lieb hat, den züchtiget; er stäupet aber einen jeglichen Sohn, den er aufnimmt. So ihr die Züchtigung erduldet, so erbeut sich euch Gott als Kindern. Das erinnert auch das 18. Kapitel.

\section[Bei der Regierung Hiskia wird Jerusalem von Sanherib]{Bei der Regierung Hiskia wird Jerusalem von Sanherib belagert}

*NehusThan) ein eherner König ist das, der die Schlange, von Gott selber in der Zeit geboten und aufgerichtet, abbricht, darum, daß sie in Mißbrauch der Abgötterei geraten war, und betrachtet darf predigen lassen, es sei NehusThan, das ist, ein ehernes Schlanglein, ein geringes Stück Erz, ein klein Kupferling, was sollte das für ein Gott sein?

\section{Summarium}

Hie gehet an die schöne Historia von dem frommen und gottseligen Könige Hiskia, welchen alle Regenten, als ein sonderlich Fürbild, mit Fleiß ansehen und seinem Exempel folgen sollen, wo sie anders Glück und Heil wollen haben. Dieser fromme König wird gerühmt aufs erste, daß er die ärgerlichen Kirchen abgethan und dem unrechten Gottesdienst mit Ernst gewehret habe. Dazu der ehernen Schlangen nicht verschonet, weil sie in einen solchen Mißbrauch kommen war, daß die Leute ihr opferten, V. 4. Zum andern, daß er Gott vertraut, an seinem Wort fest gehalten und von Gott und seinem Geboten nicht sei abgewichen, V. 5. Darum ist der HErr mit allen Gnaden und reichem Segen bei ihm gewesen, daß er von seinen Feinden wunderbarlich erlöset ist worden. Dagegen die andern, wie wir jetzt in vielen Historien
\quellseite{574}
gehöret, um ihrer Abgötterei willen in alles Unglück kommen sind.

b Rabſake heißt auf deutsch ein Erschenkne.

\section{Summarium}

Ist eine Historia, in welcher man siehet, wie der fromme König Hiskiä in großer Angst und Noth ist gewesen, daß Sanherib in Juda her-

34*
\quellseite{575}
auf zeucht und nimmt alle feste Städte ein und sendet seine Botschaft mit großer Macht für Jerusalem, dadurch er dem König Hiskia allen Trost, beide von Gott und den Menschen, ausreden will; sonderlich aber arbeitet er dahin, daß er ihm das Volk in der Stadt verzagt und abwendig mache.

In solcher Noth werden wir hören, wie der fromme König, sich an Gottes Wort gehalten und sich und seine Leute auf Gottes Hülfe vertröstet habe. Wie denn solches der gewisse und beste Trost ist, dahin wir in aller Anfechtung sehen sollen.

\section{Votum}

Helfe GOtt, daß wir auch in wahrem Glauben an ihm halten, in aller Noth ihn anrufen, und steure der Christenheit Sanherib, den Türken, und allen andern seinen und unsern Feinden, Amen.

\chapter{Das 19. Kapitel}

\section{Vorrede}

Ob die Gläubigen wohl eine Zeitlang einen harten Schweiß unter dem Kreuz und in Anfechtung leiden müssen; so haben sie doch ein Kräutlein, das heißt Heil aller Welt und heilig Wundkraut, welches ist unser lieber HErr JEsus Christus, der uns arme Heiden von allen Wunden des Teufels und Todes, die sie uns durch die Sünde gebissen, geheilet, und das rechte Heil aller Welt worden ist, und der noch hilft in geistlicher und leiblicher Noth. Denn endlich finden sich doch die liebe Angelica und die Himmelschlüssel, das ist, die Gläubigen werden zuletzt aus aller Angst gerissen, alle Thränen sollen ihnen von ihrem Angesicht abgewischt und sollen mit allen Engeln mit himmlischer Freude und ewiger Seligkeit ergötzet werden. Derhalben die Christen sich nicht sollen irren lassen, ob sie oft den Gesang traurig müssen anheben; am Ende wird sichs ausweisen, wo es hinaus will. Die Welt hebt zwar ihren Gesang hoch an, darum nicht Wunder ist, ob sie ihn endlich nicht wird erschreien können. Und wenn ihr letztlich das Halleluja wird gelegt werden, so wird sie mit Schanden müssen stille schweigen. Das unterrichtet auch der heilige Prophet Esaias am 3. Kapitel V. 10. 11. und sagt: Prediget von den Gerechten, daß sie es gut haben, denn sie werden die Frucht ihrer Werke essen. Wehe aber den Gottlosen, denn sie sind boshaftig und es wird ihnen vergolten werden, wie sie es verdienen. Davon höret auch an das 19. Kapitel.

\section[Sanheribs Trot[z] und Macht wird auf Hiskias Gebet]{Sanheribs Trot[z] und Macht wird auf Hiskias Gebet gebrochen und getrennet}
\quellseite{576}

\section{Summarium}

Dieß ist ein sehr schön und tröstlich Exempel, wie der fromme König Hiskiä in seiner größten Noth zu GOtt seine Zuflucht hat, sich vor ihm von Herzen demüthiget und um Hülfe bittet im Hause des HErrn, V. 1. Item, er sendet zu dem Propheten Jesaia, daß er auch wolle zu dem HErrn um Hülfe schreien, wider das Lästern des Königs zu Assyrien, V. 15. Solches zeiget aber an, wie ein herzliches Vertrauen der fromme König auf Gottes Güte gehabt habe. Darum hängt er mit festem Glauben am Trost und an der Zusage, die ihm der Prophet Jesaia thut.

Und dies ist der höchste und beste Gottesdienst, in Nöthen GOtt also anrufen, und sich seiner Barmherzigkeit und seines Wortes trösten, und nicht auf Menschen Gewalt, Geld noch Gunst sehen.
\quellseite{577}

\section{Summarium}

Ist eine Trost-Predigt, in welcher der Prophet Jesaia den frommen König Hiskia ermahnet, daß er auf Gottes Hülfe warten und in der Noth nicht verzagen soll. Sie ist aber sonderlich zu merken, daß Gott das Gebet derer, die auf seine Hülfe warten, hören und zu rechter Zeit helfen will. Wiederum will er den Stolz und Hochmuth der Gottlosen nicht ungestraft lassen abgehen, wie man an Sanherib und seinem Heer siehet. Sonderlich aber ist hie zu lernen, daß man in seinem Unglück nicht verzweifeln soll, sondern hoffen, Gott könne und wolle Hülfe schaffen, wo es uns gleich unmöglich dünkt. Wie denn der König Sanherib hie abgetrieben und in einer Nacht hundert und fünf und achtzig tausend Mann durch den Engel erwürget worden, und der fromme König sammt den Seinen wunderbahrlicher Weise errettet wird, V. 35. Uns zum Trost, daß wir auch lernen also auf Gottes Hülfe sehen und warten.

\section{Votum}

Der barmherzige gütige GOtt vernehme das Schreien aller, die in ihrem Elende ihn anrufen, und erhöre auch unser Gebet in unserer Noth, tröste uns und alle Betrübte durch sein Wort, und wende ab alle Gefahr, um der Vorbitte und Verdienstes unsers HErrn JEsu Christi willen, Amen.

\chapter{Das 20. Kapitel}

\section{Vorrede}

Der heilige Prophet David sagt im 41. Psalm V. 1—4.: Wohl dem, der sich des Dürftigen annimmt, den wird der HErr erretten zur bösen Zeit, der HErr wird ihn bewahren und beim Leben erhalten, und ihm lassen wohlgehen auf Erden, und nicht geben in seiner Feinde Willen. Der HErr wird ihn erquicken auf seinem Siechbette; du hilfst ihm von aller seiner Krankheit. Denn unser Glück und Unglück, Freude und Leid, Gesundheit und Krankheit, Wohlfahrt und Uebelfahrt, Leben und Tod steht nicht in des Teufels, nicht in der Welt, nicht in unserer Gewalt, sondern allein in Gottes Händen, Macht und Gewalt. Und Salomon in Sprüchen im 16. Kapitel V. 5. sagt: Ein stolzes Herz ist dem HErrn ein Greuel, und wird nicht ungestraft bleiben. Denn wie Gott die hoffärtigen Engel aus dem Himmel, und die hoffärtigen Menschen, Adam und Eva, aus dem Paradies verstoßen hat; also will er noch auf Erden die Hoffart strafen. Beide angezogene Sprüche sind erfüllet worden am Könige Hiskia, wie solches bezeuget das 20. Kapitel.

\section{Hiskias Krankheit, Lebens-Verlängerung, Ehrgeiz und Abziehen}
\quellseite{578}

\section{Summarium}

Ist endlich ein Exempel, wie Gott das Gebet der Gottseligen, so auf seine Güte hoffen, erhöret und ihnen hilft in Krankheiten und allerlei Nöthen. Sonderlich ist hie zu merken, daß es nicht Sünde ist, bitten um Gesundheit und langes Leben, und der Arznei pflegen, V. 7. So fern, daß man Gott zuvor durchs Gebet ersuche. Zum andern ist ein Exempel hie, wie es der fromme Hiskia nach seiner großen Krankheit auch verstehet, und sich versündiget mit Hoffart, darum, daß der König zu Babel seine Botschaft und Geschenke zu ihm sendet, daß er wieder gesund war worden, V. 12. Solche Botschaft und Ehre machet den König Hiskia fröhlich, daß er sich überhub und zeigte des Königs von Babel Gesandten alle seine Schätze, V. 13. Darum dräuet ihm Gott durch Jesaiam, daß alles, was er und seine Väter gesammlet haben, soll gen Babel geführet werden, und seine Nachkommen sollen des Königs von Babel Gefangene sein, V. 17. Solche Strafe nimmt Hiskia geduldig an, und erkennet, daß ers verdienet habe, V. 19.

\section{Votum}

Der treue GOtt, unser lieber Vater, erhöre auch unser Gebet, und helfe uns und allen, die seiner Hülfe bedürfen an Leib und Seele, und regiere uns, daß wir in seiner Furcht beständig leben, für Hoffart und andern Sünden uns hüten, Amen.

\chapter{Das 21. Kapitel}

\section{Vorrede}

St. Paulus sagt Röm. 2, 13.: Es ist kein Ansehen der Person vor Gott. Welche ohne Gesetz gesündiget haben, die werden auch ohne Gesetz verloren werden, und welche am Gesetze gesündiget haben, die werden durch das Gesetz verurtheilet werden; sintemal für Gott nicht, die das Gesetz hören, gerecht sind, sondern die das Gesetz thun, werden gerecht sein. Dies wiederholet er an die Colosser im 3. Capitel V. 25., da er sagt: Wer Unrecht thut, wird empfahen, was er Unrecht gethan hat, und gilt kein Ansehen der Person. Hat nun Gott die Sünde gestraft an hohen Potentaten, Königen, Fürsten und Herren, so wird er sie an gemeinen Leuten, Bürgern und Bauern freilich nicht ungestraft lassen. Derwegen wollen wir uns an andern spiegeln, wollen aus fremder Leute Schaden klug, und aus ihren Strafen fromm werden, wollen Buße thun, für Sünden uns hüten und in Gottesfurcht leben, daß wir nicht einen zornigen, sondern gnädigen Gott allzeit haben und behalten. O HErr Jesu Christe, das hilf uns allen, Amen. Darauf höret ferner an das 21. Kapitel.

\section{Historie von Manasse und Amon, den Königen Juda}
\quellseite{579}

\section{Summarium}

Ist ein schrecklich Exempel, wie der fromme König Hiskia einen ungerathenen Sohn hat, der allen Jammer und Abgötterei und Tyrannisiren wider die frommen Propheten anrichtet, V. 2. Aber hie sehet man wiederum, wie gereu-
\quellseite{580}
lich Gott solche Sünde träuet zu strafen, daß er sein Volk, gleich wie zuvor die Heiden, vertilget und in ein fremdes Land wolle lassen gefangen wegführen, V. 11. Am Ende ist eine historia, wie Amon, Manasse Sohn, König worden und von den Seinen erschlagen sei, V. 23. Aber die Aufrührer, so solchen Todschlag begangen hatten, werden darum gestraft, V. 24., nach dem Spruch: Wer das Schwert nimmt, soll durchs Schwert umkommen.

\section{Votum}

Regiere der treue liebe Gott die Jugend, daß sie in ihrer christlichen Eltern Fußstapfen treten, fromm sein, Gott ehren und ihrem Nächsten dienen; und behüte für Aufruhr, Krieg und allem andern Unglück und Sündestrafen, Amen.

\chapter{Das 22. Kapitel}

\section{Vorrede}

Wie der König Josia ihm das Gesetz-Buch Mose hat lesen lassen: also hat es im Judenthum eine sonderliche Ordnung gehabt, daß die Bücher Mose alle Sabbather für der Gemeinde sind gelesen worden, Apost. 13, 14., c. 18, 4. Dergleichen das Exempel Christi zu Nazareth, Luc. 4, 16., zeiget auch an, daß man in den Synagogen der Propheten Bücher gelesen und die Juden davon Unterredung gehalten haben. Und der Apostel St. Paulus vermahnet seine Colosser im 3. Kapitel V. 16.: Daß sie das Wort Christi reichlich unter ihnen sollen wohnen lassen in aller Weisheit, daß ein jeglicher sich selbst dadurch lehren und vermahnen möge: also wird uns auch das Alte und Neue Testament fürgelesen. Denn dieweil der Anfechtungen sehr viel sind, so ohne Unterlaß vom Teufel, der Welt und unserm eigenen Fleisch erreget werden, in Sünden und Unglauben zu stürzen, ists unmöglich, daß man sich dagegen aufhalten, wehren und trösten möge, denn allein durch das Wort Gottes. Derwegen, ach gebraucht ja fleißig das Vorlesen der heiligen Biblia zur Besserung eures Lebens, zur Stärkung eures Glaubens und zu eurem Trost in Noth und Tod, Leben und Sterben. Ferner höret nun an das 22. Kapitel.

\section{Unter Josia wird das Gesetzbuch gefunden}
\quellseite{581}

\section{Summarium}

Ist eine schöne Historia von dem frommen König Josia, wie er sich erstlich der Kirchen annimmt, V. 3., auf daß der rechte Gottesdienst gefördert werde, und wie ohngefähr das Gesetzbuch im Tempel gefunden wird, V. 8., läßt es ihm lesen und richtet alles wieder an, was von Gott geboten, aber unter der gottlosen Könige Manasse und Amons Regiment gefallen war. Daraus sehen wir, alsobald Gottes Wort von den Leuten und von der Kanzel kommt, ist’s unmöglich, daß man sich der Abgötterei und anderer Sünden, so daraus folgen, könne wehren. Sonderlich aber ist hie zu merken, daß Gott die künftige und längst verdiente Strafe gnädig darum aufhält, daß der fromme König Gottes Wort annimmt, Buße thut und sich bessert.

\section{Votum}

Wir danken GOtt für alle christliche Obrigkeit; der wolle ihm ferner auch unsere Regenten lassen befohlen sein, von ihnen und uns die verdienten Strafen abwenden, oder ja lindern, und geben, daß die heilige Biblia auch in dieser Kirche, bis an den jüngsten Tag, fleißig und fruchtbarlich vorgelesen werde, Amen.

\chapter{Das 23. Kapitel}

\section{Vorrede}

Fromme Obrigkeit, das reine und lautere Wort Gottes und christliche Kirch-Ordnung, sind große und heilsame Wohlthaten Gottes, und aller Land und Leute schönste Zierde und nützlichste Kleinod. Für diese Gnaden-Werke und Wohlthaten laßt uns heute abermal Gott dem HErrn herzlich danken und fleißig beten, daß er ihm unsere liebe Obrigkeit ferner lasse befohlen sein, auf daß wir unter ihrer christlichen Regierung Gottes Wort hören, die Gottesdienste im Glauben und mit gottseligem Leben üben mögen. Denn die Regenten sind wie die Laternen: Wie dieselbigen das Licht verwahren, daß es nicht ausgehet: also müssen nach Gott die Regenten fürnehmlich das Licht, den HErrn Christum, tragen, seine Pfleger, Nährer und Säugammen sein, soll die reine Lehre und christliche Kirchordnungen, wider alle Sturmwinde des leidigen Teufels und der argen Welt beschützt und erhalten werden. Deß sollen wir uns auch erinnern bei dem 23. Capitel.

\section[Josia reiniget den Gottesdienst. Regierung der Könige Joahas]{Josia reiniget den Gottesdienst. Regierung der Könige Joahas und Josiam}
\quellseite{582}

a Samarim) das waren sonderliche geistliche Leute, wie jetzt die allernächtigsten und strengsten Mönche sein wollen. Darum haben sie auch einen Namen Samarim, der lautet als von thätiger großer Andacht. Und das Räuchern galt bei ihnen, als ob mans der Mönche Singen und Beten in der Kirche. Denn Räucherwerk bedeutet allenthalben Gebet in der Schrift. Aber wie dieser Gebet ist, so war jener Räucherei, beide Menschenwerk ohne Gottes Geist und Wort.

b Maschith) heißt Verderbung; und man hält, es sei der Oelberg gewesen, da man die Uebelthäter hat abgethan, den wir den Galgenberg oder Rabenstein heißen.

c Grabmal) was aber solche Klammern oder Grabhügel und Grabzeichen sind, wirst du in Jeremia und Hesekiel wohl finden.

\section{Summarium}

Dies ist eine treffliche Historia, wie der fromme König Josia endlich die Predigt wieder anrichtet, daß man das Gesetzbuch öffentlich im Tempel lieset, V. 2., und hernach alle Abgötterei, Altar, Kirchen und Kirchen-Geräthe, so man zur Abgötterei brauchet, vertilget, und die abgöttischen Priester erwürget, V. 4.

Solcher Ernst gefällt Gott wohl; darum wäre es wohl zu wünschen, daß die weltliche Obrigkeit
\quellseite{583}
 dergleichen mit dem gottlosen Papstthum auch handelte, auf daß doch Gottes Zorn abgewendet und nicht der Abgötterei wiederum die Thür geöffnet würde. Weil man Klöster, Altäre, abgöttische Bilder, Bücher und dergleichen noch bleiben läßt, und nicht weg thut, wie sie werth sind. Sintemal sie zu nichts gewisser, denn zum Aergerniß dienen.

\section{Summarium}

Ist eine Historia, wie der König Josia neben dem, daß er alle Greuel und Abgötterei hat ausgefegt, auch das Osterfest wieder angerichtet habe, welches seit der Richter Zeit nie so gehalten worden. Denn weil die Leute eine äußerliche Uebung und Gottesdienst müssen haben, soll man doch in allen Dingen sonderlich auf das Wort sehen und nach demselben solche äußerliche Uebung anstellen. Solchem Exempel folgen die Söhne Josias nicht, derhalben kommt alles Unglück über sie, bis endlich das ganze Land um der Sünden willen verwüstet wird, V. 32.

\section{Votum}

Dem treuen GOtt sei Lob, Ehre, Preis und Dank, der uns von der Abgötterei auch erlöset hat; der regiere und behüte ferner Obrigkeit und Unterthanen, daß sie sich auch in der Lernung göttlichen Worts und in Gebrauch
\quellseite{584}
 christlicher Kirchordnung erzeigen, als Schülern und Dienern Christi, unsers HErrn, gebühret, Amen.

\chapter{Das 24. Kapitel}

\section{Vorrede}

Viele Sünden werden belohnet mit viel Strafen; und auf große Sünden folgen auch große Strafen; und, je mehr man sündiget an einem Ort, je näher die Strafen herzu kommen; und, je frecher und sicherer man ist im Sündenleben, je plötzlicher die Strafe herein bricht. Derwegen, ach heute ist es Zeit, daß wir alle Buße thun, und Gott um Vergebung der Sünden und Linderung der Sünde-Strafen bitten, im Namen unsers HErrn Jesu Christi. Heute ist es Zeit, daß wir von allen öffentlichen Sünden ablassen; heute ist es auch Zeit, daß man anfange, christlicher, ehrbarer, züchtiger und aufrichtiger zu leben, denn leider von vielen bisher geschehen ist. Thun wir dieses alles nicht, so folget wahrhaftig Jammer und Herzeleid mit den größten Strafen. Dieses lehret auch das 24. Kapitel.

\section[Von dreien Königen in Juda: Joakim, Jojachim und]{Von dreien Königen in Juda: Joakim, Jojachim und Zidekia}
\quellseite{585}

\section{Summarium}

Dies ist eine Historia, wie sich die Sünde und das gottlose Leben gehäufet und die Strafe je länger je näher kommen sei. Und meldet der Text fürnehmlich die Sünde Manasse, welcher Jerusalem mit unschuldigem Blut hatte voll gemacht, V. 3. Solche Sünde (spricht Gott) wolle er nicht vergeben. Darum ist gewiß zu besorgen, daß das unschuldige Blut, welches die gottlosen Papisten zu unsern Zeiten allenthalben haben vergossen, und noch, werde das Deutschland auch drücken, daß es endlich müsse zu Boden gehen.

\section{Votum}

Der treue Gott und barmherzige liebe Vater sei uns allen gnädig, lindere die wohlverdienten Strafen, die mit gottlosem Leben und unschuldigem Blutvergießen verdienet sind, und segne und behüte unsere Stadt und das ganze liebe Deutschland, Amen.

\chapter{Das 25. Kapitel}

\section{Vorrede}

Gleiche Sünden werden von dem gerechten Gott belohnet mit gleichen Strafen. Nun gehen jetzt überall nicht weniger und geringere Sünden in vollem Schwange, als im jüdischen Lande, da dies vom König NebucadNezar und hernach von Tito ist verwüstet worden. Wehe bewegen der ganzen Welt, wehe allen Land und Leuten, so man nicht heute Buße thut und in wahrer Buße auch lebet. Ach so laßt uns Gott herzlich bitten, er wolle uns unsere Sünden vergeben und die verdienten Strafen lindern und aufhalten um Jesu Christi, seines Sohnes, Verdienstes und Vorbitte willen; und ein jedes wandele und lebe nun in Seligkeit und Gerechtigkeit, in Kraft des heil. Geistes, auch wie es sein Stand und Beruf erfordert. Denn mit gläubigem Gebet löschet man Gottes Zorn, und mit Besserung des Lebens wendet man ab die Strafen oder man erlanget doch Linderung. Davon höret auch das 25. und letzte Kapitel.

\section[Von der Zerstörung der Stadt Jerusalem, und der]{Von der Zerstörung der Stadt Jerusalem, und der babylonischen Gefangenschaft}

a Jerusalem ist zuvor Jebus, item, Salem genannt gewesen, und von Sem, dem Sohn Noah, dem Könige Melchisedek, erbauet, von welches Tode bis auf diese Verwüstung 1236 Jahr sind; der Tempel, welcher im vierten Jahr Salomonis, nach dem Auszug aus Egypten, 480 Jahr erbauet, ist bis auf diese Zeit der Verwüstung gestanden 461 Jahr. Diese Verwüstung ist geschehen nach der Welt Anfang 3394 Jahr, vor Christi Geburt 566 Jahr. Die Belagerung hat anderthalb Jahr gewähret. Der Tempel ist hernachmals, nämlich nach dieser Verwüstung 123 Jahr, im 6. Jahr Darii, am dritten Tage Marti wieder vollkommen erbauet, Esra 6., vor Christi Geburt 445. In wieder sammt der Stadt Jerusalem gefangen 518 Jahr, und im andern Jahr des Kaisers Vespasiani im April, am österlichen Fest den 14. April belagert, der Tempel erobert und zum andernmal eben im Monat und auf den achten Tag des Augustmonats, wie zuvor vom Nabuchodonosor, also zum letzten von Tito, des Vespasiani Sohn, verbrannt. Die Stadt aber am achten Tage des Septembers eingenommen, geschleift und verbrennet. Dies ist geschehen nach Christi Himmelfahrt 40 Jahr, nach seiner Geburt 74 Jahr, nach dieser letzten Verwüstung 60 Jahr. Unter dem Kaiser Adriano erregten die Juden einen Aufruhr wider den Kaiser, und wollten Jerusalem und den Tempel wieder erbauen, wurden darüber fünf mal hundert und achtzig tausend Juden auf einen Tag erschlagen, Jerusalem wieder geschleift, 50 Städte und Flecken, 50 Schlösser zerstöret und verbrannt. Erobret die Verwüstung nun 1500 Jahr, und wird also ewig bleiben, wie Daniel geweissaget.
\quellseite{586}

\section{Summarium}

Ist eine schreckliche Historia, wie Gott das Königreich Juda um der Sünde willen in der Feinde Hände gegeben, verheeren und sein Volk gefangen habe lassen wegführen, V. 8. Zum Exempel uns allen, daß wir Gottes Wort folgen, in der Furcht Gottes leben, und alle Abgötterei, sammt falscher Lehre mit höchstem Ernst fliehen sollen. Denn weil Gott seines eigenen Volkes nicht verschonet hat, dem er allein, und sonst keinem andern auf Erden zur selbigen Zeit, sein Wort und rechten Gottesdienst vertrauet hatte, so dürfen wir anders nicht gedenken, er werde dergleichen Sünde auch an uns mit Krieg, Aufruhr und dergleichen Plagen, sonderlich aber mit allerlei Lügen und Irrthum strafen, wo wir uns nicht bessern.

\section{Votum}

Der barmherzige GOtt sei uns gnädig, regiere und helfe uns, daß wir uns alle bessern; und behüte uns für Krieg, Aufruhr, Irrthum und für andern zeitlichen, und dem auch ewigen, Verderben und Schaden, Amen.
